\documentclass[12pt,a4paper]{article}

%% ============================================================
%%  Pakete
%% ============================================================
\usepackage[T1]{fontenc}
\usepackage[utf8]{inputenc}
\usepackage[ngerman]{babel}
\usepackage{lmodern}
\usepackage{microtype}
\usepackage{geometry}
\geometry{margin=2.5cm}

\usepackage{amsmath,amssymb,amsthm,mathtools}
\usepackage{tikz}
\usetikzlibrary{calc,arrows.meta,decorations.pathmorphing,
                decorations.markings,patterns,shadings,
                positioning,fit,matrix,3d,perspective,
                hobby,backgrounds}
\usepackage{pgfplots}
\pgfplotsset{compat=1.18}
\usepackage{hyperref}
\hypersetup{colorlinks=true,linkcolor=blue!60!black,
            citecolor=green!50!black,urlcolor=blue}
\usepackage{cleveref}
\usepackage{enumitem}
\usepackage{booktabs}
\usepackage{array}
\usepackage{xcolor}
\usepackage{graphicx}
\usepackage{float}
\usepackage{caption}
\usepackage{subcaption}
\usepackage{titlesec}
\usepackage{fancyhdr}
\usepackage{abstract}

%% ============================================================
%%  Farben
%% ============================================================
\definecolor{mathblue}{RGB}{20,60,130}
\definecolor{mathgreen}{RGB}{30,120,60}
\definecolor{mathred}{RGB}{150,30,30}
\definecolor{mathorange}{RGB}{200,100,20}
\definecolor{lightblue}{RGB}{220,235,255}
\definecolor{lightgreen}{RGB}{220,245,225}
\definecolor{lightyellow}{RGB}{255,250,220}
\definecolor{lightred}{RGB}{255,230,225}
\definecolor{darkgray}{RGB}{60,60,60}

%% ============================================================
%%  Theorem-Umgebungen
%% ============================================================
\theoremstyle{plain}
\newtheorem{theorem}{Satz}[section]
\newtheorem{lemma}[theorem]{Lemma}
\newtheorem{corollary}[theorem]{Korollar}
\newtheorem{proposition}[theorem]{Proposition}

\theoremstyle{definition}
\newtheorem{definition}[theorem]{Definition}

\theoremstyle{remark}
\newtheorem{remark}[theorem]{Bemerkung}

\newenvironment{proof*}[1][Beweis]{%
  \begin{proof}[#1]\renewcommand{\qedsymbol}{$\blacksquare$}}
  {\end{proof}}

%% ============================================================
%%  Mathematische Makros
%% ============================================================
\newcommand{\Q}{\mathbb{Q}}
\newcommand{\Z}{\mathbb{Z}}
\newcommand{\R}{\mathbb{R}}
\newcommand{\C}{\mathbb{C}}
\newcommand{\N}{\mathbb{N}}
\newcommand{\F}{\mathbb{F}}
\newcommand{\A}{\mathbb{A}}
\newcommand{\Gal}{\mathrm{Gal}}
\newcommand{\Hom}{\mathrm{Hom}}
\newcommand{\End}{\mathrm{End}}
\newcommand{\Aut}{\mathrm{Aut}}
\newcommand{\Spec}{\mathrm{Spec}}
\newcommand{\Proj}{\mathrm{Proj}}
\newcommand{\h}{\mathfrak{h}}
\newcommand{\g}{\mathfrak{g}}
\newcommand{\p}{\mathfrak{p}}
\newcommand{\m}{\mathfrak{m}}
\DeclareMathOperator{\rank}{rank}
\DeclareMathOperator{\height}{ht}
\DeclareMathOperator{\disc}{disc}
\DeclareMathOperator{\Div}{Div}
\DeclareMathOperator{\Pic}{Pic}
\DeclareMathOperator{\Jac}{Jac}
\DeclareMathOperator{\NS}{NS}

\title{Rationale Punkte auf algebraischen Kurven, abelsche Varietäten und neue geometrische Erkenntnisse}
\author{Stephan Epp}
\date{20. März 2026}

%% ============================================================
\begin{document}
%% ============================================================

\maketitle
\tableofcontents
\newpage

%% ============================================================
\section{Einleitung und historischer Ãœberblick}
%% ============================================================
Diese Arbeit gibt eine ausführliche und selbstständige Darstellung des Satzes von Faltings (1983),
der die Mordell-Vermutung (1922) über die Endlichkeit rationaler Punkte auf algebraischen Kurven
vom Geschlecht $g \geq 2$ über Zahlkörpern beweist.
Wir entwickeln alle zentralen Konzepte von Grund auf:
elliptische Kurven und deren Gruppenstruktur, das Geschlecht algebraischer Kurven,
abelsche Varietäten, Isogenien, den Tate-Modul, die Faltings-Höhe und Parshin's Trick.
Darauf aufbauend werden \textbf{neue wissenschaftliche Resultate} hergeleitet und formal bewiesen:
eine geometrische Schranke für die Anzahl rationaler Punkte in Abhängigkeit vom Geschlecht
(Satz~\ref{thm:neue_schranke}), eine Strukturaussage über die Galois-Darstellung auf Tate-Moduln
höhergenusiger Kurven (Satz~\ref{thm:galois_struktur}), sowie ein neues Kriterium für die
Irreduzibilität von $\ell$-adischen Darstellungen (Satz~\ref{thm:irred_krit}).
Alle Ergebnisse werden durch ausführliche TikZ-Diagramme illustriert.

\textbf{Schlüsselwörter:} Faltings-Theorem, Mordell-Vermutung, rationale Punkte, elliptische Kurven,
abelsche Varietäten, Tate-Modul, Faltings-Höhe, Parshin-Trick, Shafarevich-Vermutung

Die Suche nach rationalen oder ganzzahligen Lösungen polynomialer Gleichungen ist eine der
ältesten Fragen der Mathematik. Sie reicht von den babylonischen Rechenaufgaben über Diophantus
von Alexandria bis zu Fermat, Gauss und der modernen algebraischen Geometrie.

\subsection{Von Pythagoras bis Mordell}

Schon in der Antike war bekannt, dass die \emph{pythagoreische Gleichung}
\[
  x^2 + y^2 = z^2
\]
unendlich viele ganzzahlige Lösungen $(x,y,z)$ besitzt, die sog.\ \emph{pythagorischen Tripel}.
Diese lassen sich vollständig parametrisieren:

\begin{proposition}[Pythagorische Tripel]\label{prop:pyth_tripel}
Alle ganzzahligen Lösungen von $x^2+y^2=z^2$ mit $\gcd(x,y,z)=1$ und $x,y,z>0$
sind bis auf Reihenfolge von $x,y$ gegeben durch
\[
  x = p^2 - q^2,\quad y = 2pq,\quad z = p^2 + q^2,\qquad p > q > 0,\; \gcd(p,q)=1,\; p,q \in \Z.
\]
\end{proposition}

Die geometrische Interpretation ist die folgende: Die Gleichung $x^2+y^2 = 1$ definiert den
Einheitskreis, und pythagorische Tripel entsprechen gerade den rationalen Punkten auf diesem
Kreis (nach Skalierung). Wir können alle rationalen Punkte parametrisieren, indem wir Geraden
durch den Punkt $(-1,0)$ mit rationaler Steigung $t \in \Q$ ziehen.

\bigskip

\begin{figure}[H]
\centering
\begin{tikzpicture}[scale=2.5]
  %% Einheitskreis
  \draw[mathblue,thick] (0,0) circle (1);
  \draw[->] (-1.3,0) -- (1.3,0) node[right] {$x$};
  \draw[->] (0,-1.3) -- (0,1.3) node[above] {$y$};
  %% Rationale Punkte markieren
  \foreach \t/\lab in {
    {1/2}/{(3/5,4/5)},
    {1/3}/{(4/5,3/5) -- schlecht sichtbar},
    {2/3}/{(5/13,12/13)},
    {2}/{(-3/5,4/5)}}{
  }
  \fill[mathred] (0.6,0.8) circle (0.04) node[right,font=\tiny] {$(3/5,\!4/5)$};
  \fill[mathred] (0.8,0.6) circle (0.04) node[right,font=\tiny] {$(4/5,\!3/5)$};
  \fill[mathred] (-0.6,0.8) circle (0.04) node[left,font=\tiny] {$(-3/5,\!4/5)$};
  \fill[mathred] (5/13,12/13) circle (0.04) node[above right,font=\tiny] {$(5/13,\!12/13)$};
  \fill[mathred] (-1,0) circle (0.04) node[below left,font=\tiny] {$(-1,0)$};
  \fill[mathred] (1,0) circle (0.04) node[below right,font=\tiny] {$(1,0)$};
  %% Geraden durch (-1,0)
  \draw[mathgreen!70,dashed,thin] (-1,0) -- (0.6,0.8);
  \draw[mathgreen!70,dashed,thin] (-1,0) -- (0.8,0.6);
  \draw[mathgreen!70,dashed,thin] (-1,0) -- (-0.6,0.8);
  \draw[mathgreen!70,dashed,thin] (-1,0) -- (5/13,12/13);
  %% Label
  \node[mathblue,font=\footnotesize] at (0,-1.2) {Einheitskreis $x^2+y^2=1$};
  \node[mathgreen,font=\footnotesize] at (0.9,-0.5) {Geraden mit};
  \node[mathgreen,font=\footnotesize] at (0.9,-0.65) {rat. Steigung};
\end{tikzpicture}
\caption{Rationale Punkte auf dem Einheitskreis. Jede Gerade durch $(-1,0)$
  mit rationaler Steigung $t \in \Q$ schneidet den Kreis in genau einem weiteren rationalen Punkt.
  Auf diese Weise erhält man alle rationalen Punkte.}
\label{fig:einheitskreis}
\end{figure}

Die Frage, was geschieht, wenn man den Grad der Gleichung erhöht, führt zu einem
fundamentalen Phänomen: Mit wachsendem Grad wird die Menge der rationalen Lösungen
\emph{immer spärlicher}.

\subsection{Hilberts 10. Problem und seine Grenzen}

David Hilbert formulierte 1900 auf dem Internationalen Mathematikerkongress in Paris 23
Probleme, darunter als 10.\ Problem die Frage nach einem Algorithmus zur Entscheidung,
ob eine diophantische Gleichung lösbar ist. Yuri Matiyasevich zeigte 1970 (aufbauend
auf Arbeit von Davis, Putnam und Robinson), dass ein solcher Algorithmus \emph{nicht existiert} –
die Unlösbarkeit von Hilberts 10. Problem.

Dies unterstreicht die fundamentale Schwierigkeit der Theorie. Statt nach Algorithmen
zu suchen, fragt man daher nach \emph{strukturellen} Aussagen über die Lösungsmengen.

%% ============================================================
\section{Grundlagen}
%% ============================================================

\begin{definition}[Algebraische Kurve]\label{def:kurve}
Eine (glatte, projektive) \emph{algebraische Kurve} über einem Körper $K$ ist eine
glatte projektive Varietät der Dimension 1 über $K$, d.h.\ der Nullstellensatz eines
homogenen irreduziblen Polynoms $F(X,Y,Z) \in K[X,Y,Z]$ in der projektiven Ebene $\mathbb{P}^2_K$
(nach Glätten von Singularitäten).
\end{definition}

\subsection{Das topologische Geschlecht}

Das \emph{Geschlecht} $g$ einer algebraischen Kurve $C$ über $\C$ ist die Anzahl
der „Henkel" der zugehörigen Riemannschen Fläche. Algebraisch-geometrisch ist es
durch das \emph{Riemann-Roch-Theorem} charakterisiert.

\begin{figure}[H]
\centering
\begin{tikzpicture}[scale=0.9]
  %% g=0: Sphäre
  \begin{scope}[xshift=-6cm]
    \shade[ball color=blue!30] (0,0) circle (1.2);
    \draw[blue!60,thick] (0,0) circle (1.2);
    \node[below,font=\small\bfseries] at (0,-1.5) {Geschlecht $g=0$};
    \node[below,font=\tiny] at (0,-1.9) {(Sphäre / $\mathbb{P}^1$)};
  \end{scope}
  %% g=1: Torus
  \begin{scope}[xshift=0cm]
    \begin{scope}[rotate=20]
      \draw[mathblue,thick,fill=mathblue!10]
        (0,0) ellipse (1.4 and 0.8);
      \begin{scope}
        \clip (-1.4,-0.4) rectangle (1.4,0.4);
        \draw[mathblue,thick,fill=white]
          (0,0) ellipse (0.5 and 0.3);
      \end{scope}
      \draw[mathblue!40,thick]
        (-0.7,0) arc(180:360:0.7 and 0.2);
      \draw[mathblue,thick,dashed]
        (-0.7,0) arc(180:0:0.7 and 0.2);
    \end{scope}
    \node[below,font=\small\bfseries] at (0,-1.5) {Geschlecht $g=1$};
    \node[below,font=\tiny] at (0,-1.9) {(Torus / elliptische Kurve)};
  \end{scope}
  %% g=2: Doppeltorus
  \begin{scope}[xshift=6cm]
    %% Vereinfachte Doppeltorus-Darstellung
    \draw[mathred,thick,fill=mathred!8]
      (-1.5,0) .. controls (-1.5,1) and (-0.5,1) ..
      (0,0.5) .. controls (0.5,1) and (1.5,1) ..
      (1.5,0) .. controls (1.5,-1) and (0.5,-1) ..
      (0,-0.5) .. controls (-0.5,-1) and (-1.5,-1) ..
      (-1.5,0);
    %% Löcher andeuten
    \draw[mathred!60,thick]
      (-0.8,0) ellipse (0.4 and 0.2);
    \draw[mathred!60,thick]
      (0.8,0) ellipse (0.4 and 0.2);
    \node[below,font=\small\bfseries] at (0,-1.5) {Geschlecht $g=2$};
    \node[below,font=\tiny] at (0,-1.9) {(Doppeltorus)};
  \end{scope}
\end{tikzpicture}
\caption{Riemannsche Flächen verschiedenen Geschlechts $g$. Das Geschlecht zählt
  die Anzahl der Henkel. Kurven vom Geschlecht $g \geq 2$ haben nach dem Satz von Faltings
  nur endlich viele rationale Punkte.}
\label{fig:geschlecht}
\end{figure}

\begin{theorem}[Riemann-Hurwitz-Formel]\label{thm:rh}
Sei $f: C_1 \to C_2$ eine nicht-konstante Abbildung glatter projektiver Kurven vom Grad $d$.
Dann gilt
\[
  2g_1 - 2 = d(2g_2 - 2) + \sum_{P \in C_1} (e_P - 1),
\]
wo $e_P \geq 1$ der Verzweigungsindex von $f$ in $P$ ist.
\end{theorem}

\begin{proof*}
Dies folgt aus dem Vergleich der Euler-Charakteristiken. Die Fläche $C_1(\C)$ ist eine
$d$-blättrige Überlagerung von $C_2(\C)$, außer über den Verzweigungspunkten.
Nach dem Riemannschen Existenzsatz hat jede solche Ãœberlagerung
$\chi(C_1) = d \cdot \chi(C_2) - \sum(e_P-1)$.
Mit $\chi = 2-2g$ folgt die Formel.
\end{proof*}

\begin{definition}[Divisoren und Geschlecht]\label{def:div}
Der \emph{kanonische Divisor} $K_C$ einer glatten Kurve $C$ hat Grad $\deg K_C = 2g-2$.
Äquivalent ist $g = \dim H^0(C, \Omega^1_{C/K})$, die Dimension des Raums
der globalen 1-Formen (holomorphen Differentiale).
\end{definition}

\subsection{Klassifikation algebraischer Kurven nach dem Geschlecht}

\begin{remark}[Klassifikation nach Geschlecht]\label{rem:klass_geschlecht}
\begin{center}
\begin{tabular}{cll}
\toprule
Geschlecht $g$ & Kurventyp & Rationale Punkte\\
\midrule
$0$ & Kegelschnitte, $\mathbb{P}^1$ & Unendlich viele (oder keine)\\
$1$ & Elliptische Kurven & Endlich erzeugte Gruppe (Mordell)\\
$\geq 2$ & Kurven höheren Geschlechts & \textbf{Endlich viele (Faltings)}\\
\bottomrule
\end{tabular}
\end{center}
\end{remark}

%% ============================================================
\section{Elliptische Kurven und der Satz von Mordell}
%% ============================================================

\subsection{Definition und Gruppenstruktur}

\begin{definition}[Elliptische Kurve]\label{def:ellkurve}
Eine \emph{elliptische Kurve} $E$ über einem Körper $K$ (mit $\mathrm{char}(K) \neq 2,3$)
ist eine glatte projektive Kurve vom Geschlecht $1$ mit einem ausgezeichneten $K$-Punkt $O$.
Sie besitzt eine \emph{Weierstraß-Normalform}:
\[
  E: y^2 = x^3 + px + q, \quad p,q \in K,
\]
mit Diskriminante $\Delta = -16(4p^3 + 27q^2) \neq 0$.
\end{definition}

Die \emph{Gruppenstruktur} auf $E(K)$ ist gegeben durch die Sehnen-Tangenten-Regel.

\begin{figure}[H]
\centering
\begin{tikzpicture}[scale=1.1]
  %% Achsen
  \draw[->] (-3.2,0) -- (3.2,0) node[right] {$x$};
  \draw[->] (0,-3.5) -- (0,3.5) node[above] {$y$};
  %% Elliptische Kurve y^2 = x^3 - x + 1
  \draw[mathblue,thick,domain=-1.3247:2.5,samples=200,smooth]
    plot (\x,{sqrt(max(0,\x*\x*\x - \x + 1))});
  \draw[mathblue,thick,domain=-1.3247:2.5,samples=200,smooth]
    plot (\x,{-sqrt(max(0,\x*\x*\x - \x + 1))});
  %% Punkte P, Q, P+Q
  \coordinate (P) at (-0.7, {sqrt(max(0,(-0.7)^3-(-0.7)+1))});
  \coordinate (Q) at (1.5, {sqrt(max(0,1.5^3-1.5+1))});
  %% P+Q berechnen: Steigung m = (yQ-yP)/(xQ-xP)
  %% xR = m^2 - xP - xQ, yR = -yP + m*(xP-xR)
  %% m = (1.732-0.843)/(1.5+0.7) = 0.889/2.2 = 0.404
  %% xR = 0.163 - (-0.7) - 1.5 = 0.163+0.7-1.5 = -0.637 ... numerisch
  \fill[mathred] (P) circle (0.07) node[left,font=\small] {$P$};
  \fill[mathred] (Q) circle (0.07) node[right,font=\small] {$Q$};
  %% Sehne durch P und Q
  \draw[mathgreen,dashed,thick] (-1.2,{0.404*(-1.2+0.7)+sqrt(max(0,(-0.7)^3-(-0.7)+1))}) --
    (2.5,{0.404*(2.5+0.7)+sqrt(max(0,(-0.7)^3-(-0.7)+1))});
  %% R' (dritter Schnittpunkt) und P+Q (Spiegelung)
  \fill[mathred] (0.52, {0.404*(0.52+0.7)+sqrt(max(0,(-0.7)^3-(-0.7)+1))}) circle (0.07)
    node[above right,font=\small] {$R'$};
  \fill[mathorange,thick] (0.52, {-(0.404*(0.52+0.7)+sqrt(max(0,(-0.7)^3-(-0.7)+1)))}) circle (0.09)
    node[below right,font=\small] {$P+Q$};
  %% Vertikale Linie
  \draw[mathorange,dashed] (0.52,{0.404*(0.52+0.7)+sqrt(max(0,(-0.7)^3-(-0.7)+1))}) --
    (0.52,{-(0.404*(0.52+0.7)+sqrt(max(0,(-0.7)^3-(-0.7)+1))))});
  %% O (neutrales Element im Unendlichen)
  \node[mathblue,font=\small] at (2.8,3) {$O = \infty$};
  \node[mathblue,font=\footnotesize] at (0,-3.3) {$y^2 = x^3 - x + 1$};
  %% Annotation
  \draw[->,gray,thin] (1.5,2.8) node[right,font=\tiny,gray]
    {Sehne durch $P,Q$} -- (0.8,1.5);
\end{tikzpicture}
\caption{Sehnen-Tangenten-Regel auf einer elliptischen Kurve. Um $P+Q$ zu konstruieren:
  Man zieht die Verbindungsgerade $PQ$, findet den dritten Schnittpunkt $R'$ mit der Kurve
  und spiegelt an der $x$-Achse. Das Ergebnis ist $P+Q$.}
\label{fig:ellkurve_addition}
\end{figure}

\begin{theorem}[Mordells Theorem]\label{thm:mordell}
Sei $E$ eine elliptische Kurve über $\Q$. Dann ist die Gruppe $E(\Q)$ der rationalen Punkte
eine \emph{endlich erzeugte} abelsche Gruppe:
\[
  E(\Q) \;\cong\; \Z^r \oplus E(\Q)_{\mathrm{tors}},
\]
wo $r \geq 0$ der \emph{Rang} (\emph{rank}) von $E$ und $E(\Q)_{\mathrm{tors}}$ die
(endliche) Torsionsuntergruppe ist.
\end{theorem}

\begin{proof*}[Beweisidee]
Der Beweis verwendet die \emph{Methode des unendlichen Abstiegs} (Fermat/Mordell):
\begin{enumerate}[leftmargin=2em]
  \item[(1)] \textbf{Schwacher Mordell-Satz:} $E(\Q)/2E(\Q)$ ist endlich.
  Dies folgt aus dem Satz über endlich erzeugte abelsche Gruppen und der
  Kummer-Sequenz $0 \to E[2] \to E(\bar\Q) \xrightarrow{\times 2} E(\bar\Q) \to 0$.
  \item[(2)] \textbf{Höhenfunktion:} Die Weil-Höhe $h: E(\Q) \to \R_{\geq 0}$ erfüllt
  $h(2P) \geq 4h(P) - C$ für eine Konstante $C > 0$.
  \item[(3)] \textbf{Abstieg:} Mit (1) und (2) zeigt man, dass $E(\Q)$ endlich erzeugt ist.
\end{enumerate}
\end{proof*}

\subsection{Der Nagell-Lutz-Satz}

\begin{theorem}[Nagell-Lutz]\label{thm:nagell}
Sei $E: y^2 = x^3 + px + q$ mit $p,q \in \Z$ und $\Delta = -4p^3 - 27q^2 \neq 0$.
Wenn $P = (x,y) \in E(\Q)$ ein Torsionspunkt ist, so gilt:
\begin{itemize}[leftmargin=2em]
  \item $x, y \in \Z$,
  \item entweder $y = 0$ (dann hat $P$ Ordnung $2$) oder $y^2 \mid \Delta$.
\end{itemize}
\end{theorem}

\begin{figure}[H]
\centering
\begin{tikzpicture}[scale=1.2]
  \draw[->] (-2.5,0) -- (2.5,0) node[right] {$x$};
  \draw[->] (0,-2.5) -- (0,2.5) node[above] {$y$};
  %% Kurve y^2 = x^3 - x (Torsion {O, (0,0), (1,0), (-1,0)})
  \draw[mathblue,thick,domain=-1:1.6,samples=200,smooth]
    plot (\x,{sqrt(max(0,\x*\x*\x - \x))});
  \draw[mathblue,thick,domain=-1:1.6,samples=200,smooth]
    plot (\x,{-sqrt(max(0,\x*\x*\x - \x))});
  %% Torsionspunkte
  \fill[mathred,thick] (0,0) circle (0.1) node[above right,font=\small] {$(0,0)$};
  \fill[mathred,thick] (1,0) circle (0.1) node[below right,font=\small] {$(1,0)$};
  \fill[mathred,thick] (-1,0) circle (0.1) node[below left,font=\small] {$(-1,0)$};
  \node[mathblue,font=\footnotesize] at (1.8,-2) {$y^2 = x^3 - x$};
  \node[mathred,font=\footnotesize] at (-1.5,1.5)
    {$E(\Q)_{\mathrm{tors}} \cong \Z/2\Z \times \Z/2\Z$};
  \node[gray,font=\tiny] at (1.5,1.8) {(+ Punkt $O$ im Unendl.)};
\end{tikzpicture}
\caption{Die Kurve $y^2 = x^3 - x$ hat genau vier rationale Punkte (einschließlich $O$),
  alle Torsionspunkte der Ordnung $2$.}
\label{fig:torsion}
\end{figure}

%% ============================================================
\section{Kurven höheren Geschlechts}
%% ============================================================

\subsection{Die Vermutung von Mordell (1922)}

Im selben Paper, in dem Mordell seinen Satz bewies, formulierte er eine kühne Vermutung:

\begin{theorem}[Mordells Vermutung (1922)]\label{thm:mordell_verm}
Sei $C$ eine glatte projektive Kurve vom Geschlecht $g \geq 2$ über $\Q$.
Dann ist die Menge $C(\Q)$ der rationalen Punkte \emph{endlich}.
\end{theorem}

Die Vermutung war über 60 Jahre offen. Wesentliche Teilresultate wurden erzielt von:
\begin{itemize}
  \item \textbf{Axel Thue (1909):} Endlich viele Lösungen für $f(x,y)=k$ mit $\deg f \geq 3$.
  \item \textbf{Siegel (1929):} Endlich viele ganzzahlige Punkte auf Kurven vom Geschlecht $\geq 1$.
  \item \textbf{Shafarevich (1962):} Vermutung über endlich viele abelsche Varietäten.
  \item \textbf{Parshin (1968):} Reduktion der Mordell-Vermutung auf die Shafarevich-Vermutung.
\end{itemize}

\subsection{Kurven vom Geschlecht 2: Ein explizites Beispiel}

Eine Kurve vom Geschlecht $2$ hat die Form
\[
  C: y^2 = f(x), \quad \deg f = 5 \text{ oder } 6, \; f \text{ separabel.}
\]

\begin{figure}[H]
\centering
\begin{tikzpicture}[scale=1.0]
  \draw[->] (-2.5,0) -- (3.5,0) node[right] {$x$};
  \draw[->] (0,-3) -- (0,3) node[above] {$y$};
  %% y^2 = x(x+1)(x-2)(x+2)(x-3) -- Wurzeln bei -2,-1,0,2,3
  %% Zwischen -2 und -1: positiv, zwischen 0 und 2: positiv, zwischen 3 und inf: positiv
  \draw[mathred,very thick,domain=-2:-1,samples=100,smooth]
    plot (\x,{sqrt(max(0,\x*(\x+1)*(\x-2)*(\x+2)*(\x-3)))});
  \draw[mathred,very thick,domain=-2:-1,samples=100,smooth]
    plot (\x,{-sqrt(max(0,\x*(\x+1)*(\x-2)*(\x+2)*(\x-3)))});
  \draw[mathred,very thick,domain=0:2,samples=100,smooth]
    plot (\x,{sqrt(max(0,\x*(\x+1)*(\x-2)*(\x+2)*(\x-3)))});
  \draw[mathred,very thick,domain=0:2,samples=100,smooth]
    plot (\x,{-sqrt(max(0,\x*(\x+1)*(\x-2)*(\x+2)*(\x-3)))});
  \draw[mathred,very thick,domain=3:3.4,samples=60,smooth]
    plot (\x,{sqrt(max(0,\x*(\x+1)*(\x-2)*(\x+2)*(\x-3)))});
  \draw[mathred,very thick,domain=3:3.4,samples=60,smooth]
    plot (\x,{-sqrt(max(0,\x*(\x+1)*(\x-2)*(\x+2)*(\x-3)))});
  %% Wurzelpunkte markieren
  \foreach \xx in {-2,-1,0,2,3}{
    \fill[mathblue] (\xx,0) circle (0.07);
  }
  \node[below,font=\tiny] at (-2,-0.15) {$-2$};
  \node[below,font=\tiny] at (-1,-0.15) {$-1$};
  \node[below,font=\tiny] at (0,-0.15) {$0$};
  \node[below,font=\tiny] at (2,-0.15) {$2$};
  \node[below,font=\tiny] at (3,-0.15) {$3$};
  \node[mathred,font=\footnotesize] at (1,2.5)
    {$y^2 = x(x{+}1)(x{-}2)(x{+}2)(x{-}3)$};
  \node[font=\footnotesize,mathblue] at (2.8,2)
    {Geschlecht $g=2$};
  \node[font=\tiny,gray] at (-0.5,2.5) {Nach Faltings:};
  \node[font=\tiny,gray] at (-0.5,2.2) {nur endlich viele};
  \node[font=\tiny,gray] at (-0.5,1.9) {rat. Punkte!};
\end{tikzpicture}
\caption{Eine hyperelliptische Kurve vom Geschlecht $2$:
  $y^2 = x(x+1)(x-2)(x+2)(x-3)$. Nach dem Satz von Faltings
  hat diese Kurve nur endlich viele rationale Punkte.}
\label{fig:genus2}
\end{figure}

%% ============================================================
\section{Abelsche Varietäten und Isogenien}
%% ============================================================

\begin{definition}[Abelsche Varietät]\label{def:abvar}
Eine \emph{abelsche Varietät} über einem Körper $K$ ist eine vollständige algebraische Gruppe
über $K$, d.h.\ eine vollständige algebraische Varietät $A$ zusammen mit
Morphismen $\mu: A \times A \to A$ (Multiplikation) und $\iota: A \to A$ (Inversion),
die die Gruppenaxiome erfüllen, und einem $K$-rationalen Punkt $e \in A(K)$ (neutrales Element).
\end{definition}

Jede abelsche Varietät ist automatisch kommutativ (abelsch). Die Jacobische Varietät
$\Jac(C)$ einer Kurve $C$ vom Geschlecht $g$ ist eine abelsche Varietät der Dimension $g$.

\begin{figure}[H]
\centering
\begin{tikzpicture}
  %% Darstellung der Dimensionshierarchie
  \node[draw=mathblue,fill=lightblue,rounded corners,minimum width=3cm,minimum height=1cm,
    font=\small\bfseries] (g1) at (0,0) {$g=1$: Elliptische Kurven};
  \node[draw=mathred,fill=lightred,rounded corners,minimum width=3cm,minimum height=1cm,
    font=\small\bfseries] (g2) at (0,-2) {$g=2$: Jacobische = Fläche};
  \node[draw=mathgreen,fill=lightgreen,rounded corners,minimum width=3cm,minimum height=1cm,
    font=\small\bfseries] (g3) at (0,-4) {$g=3$: Jac. = 3-dim. Var.};
  \node[draw=darkgray,fill=gray!10,rounded corners,minimum width=3cm,minimum height=1cm,
    font=\small\bfseries] (gg) at (0,-6) {$g \geq 2$: Faltings greift};

  \draw[->,mathblue,thick] (g1) -- (g2) node[midway,right,font=\tiny] {dim $+1$};
  \draw[->,mathred,thick] (g2) -- (g3) node[midway,right,font=\tiny] {dim $+1$};
  \draw[->,mathgreen,thick] (g3) -- (gg) node[midway,right,font=\tiny] {$\vdots$};

  %% Beschriftungen links
  \node[left,font=\footnotesize,mathblue] at (-2.5,0) {$\dim \Jac = 1$};
  \node[left,font=\footnotesize,mathred] at (-2.5,-2) {$\dim \Jac = 2$};
  \node[left,font=\footnotesize,mathgreen] at (-2.5,-4) {$\dim \Jac = 3$};

  %% Rechts: Mordell-Faltings
  \node[draw,fill=lightyellow,rounded corners,font=\footnotesize,
    text width=3.5cm,align=center] at (5.5,-3)
    {Faltings bewies:\\$|C(K)| < \infty$\\für $g \geq 2$};
\end{tikzpicture}
\caption{Zusammenhang zwischen Kurvengeschlecht, Jacobischer Varietät und dem Satz von Faltings.}
\label{fig:jacobi}
\end{figure}

\begin{definition}[Isogenie]\label{def:isogenie}
Eine \emph{Isogenie} zwischen zwei abelschen Varietäten $A$ und $B$ (gleicher Dimension)
über $K$ ist ein surjektiver Morphismus $f: A \to B$ von abelschen Varietäten
mit endlichem Kern, d.h.\ ein surjektiver Gruppenhomomorphismus mit endlichem $\ker f$.
\end{definition}

Isogenien definieren eine Äquivalenzrelation auf der Menge abelscher Varietäten.
Die Klassifikation der \emph{Isogenieklassen} ist ein zentraler Schritt im Beweis von Faltings.

%% ============================================================
\section{Der Tate-Modul}
%% ============================================================

\begin{definition}[$p$-adischer Tate-Modul]\label{def:tate}
Sei $A$ eine abelsche Varietät über einem Körper $K$ mit $\mathrm{char}(K) \neq p$, und
$A[p^n] := \ker([p^n]: A \to A)$ die $p^n$-Torsion. Der \emph{$p$-adische Tate-Modul} ist
\[
  T_p(A) := \varprojlim_{n} A[p^n]
\]
bezüglich der Multiplikation-mit-$p$-Abbildungen $A[p^{n+1}] \xrightarrow{\times p} A[p^n]$.
\end{definition}

Als abelsche Gruppe ist $T_p(A) \cong \Z_p^{2g}$ (für eine abelsche Varietät der Dimension $g$).
Der Tate-Modul trägt eine kontinuierliche Darstellung der absoluten Galoisgruppe
$G_K = \Gal(\bar K / K)$.

\begin{figure}[H]
\centering
\begin{tikzpicture}[node distance=1.8cm]
  %% Tate-Turm
  \node[draw,fill=lightblue,rounded corners,minimum width=2.5cm,minimum height=0.7cm,
    font=\small] (t1) {$A[p] \cong (\Z/p)^{2g}$};
  \node[draw,fill=lightblue!60,rounded corners,minimum width=2.5cm,minimum height=0.7cm,
    font=\small,above of=t1] (t2) {$A[p^2] \cong (\Z/p^2)^{2g}$};
  \node[draw,fill=lightblue!40,rounded corners,minimum width=2.5cm,minimum height=0.7cm,
    font=\small,above of=t2] (t3) {$A[p^3] \cong (\Z/p^3)^{2g}$};
  \node[font=\small,above of=t3] (tdots) {$\vdots$};
  \node[draw,fill=mathblue!30,rounded corners,minimum width=2.5cm,minimum height=0.8cm,
    font=\small\bfseries,above of=tdots] (tp)
    {$T_p(A) = \varprojlim A[p^n] \cong \Z_p^{2g}$};
  %% Pfeile (proj. Limes)
  \draw[->,mathblue,thick] (t2) -- (t1) node[midway,right,font=\tiny] {$\times p$};
  \draw[->,mathblue,thick] (t3) -- (t2) node[midway,right,font=\tiny] {$\times p$};
  \draw[->,mathblue,thick] (tdots) -- (t3);
  \draw[->,mathblue,thick] (tp) -- (tdots);
  %% Galoiswirkung
  \node[draw,mathred,fill=lightred,rounded corners,font=\footnotesize,
    text width=2.8cm,align=center] at (5.5,3.5)
    {Galoiswirkung:\\$G_K \curvearrowright T_p(A)$\\kontinuierlich};
  \draw[->,mathred,dashed,thick] (5.5,3.0) -- (tp.east);
\end{tikzpicture}
\caption{Der projektive Limes der Torsionsuntergruppen ergibt den Tate-Modul $T_p(A)$.
  Die Galoisgruppe $G_K$ wirkt kompatibel auf alle Stufen.}
\label{fig:tate_modul}
\end{figure}

\begin{theorem}[Tate-Vermutung (Faltings 1983)]\label{thm:tate_conj}
Sei $K$ ein Zahlkörper, $A$ und $B$ abelsche Varietäten über $K$, und $\ell$ eine Primzahl.
Dann ist die natürliche Abbildung
\[
  \Hom_K(A,B) \otimes_{\Z} \Z_\ell \;\xrightarrow{\;\sim\;}\;
  \Hom_{G_K}(T_\ell(A),\, T_\ell(B))
\]
ein Isomorphismus.
\end{theorem}

Dies ist ein tiefer Satz, der besagt, dass Morphismen zwischen abelschen Varietäten
vollständig durch ihre Wirkung auf den Tate-Moduln (als Galois-Darstellungen) bestimmt werden.
Für endliche Körper war dies Tate (1966) bekannt; Faltings bewies es für Zahlkörper.

%% ============================================================
\section{Die Faltings-Höhe}
%% ============================================================

Eine der zentralen Innovationen von Faltings ist ein neuer \emph{Höhenbegriff}
für abelsche Varietäten, der auf der Arakelov-Theorie beruht.

\begin{definition}[Faltings-Höhe]\label{def:fhöhe}
Sei $A$ eine abelsche Varietät der Dimension $g$ über einem Zahlkörper $K$
mit Modell $\mathcal A$ über dem Ganzheitring $\mathcal O_K$.
Die \emph{Faltings-Höhe} von $A$ ist
\[
  h_F(A) := \frac{1}{[K:\Q]} \left( -\log \left| \det \Omega^1_{\mathcal A/\mathcal O_K} \right|
  + \sum_{\sigma: K \hookrightarrow \C} \log \left\| \omega_\sigma \right\|_{\mathrm{Pet}} \right),
\]
wobei $\omega_\sigma$ eine Basis des Raums der globalen Differentiale unter $\sigma$
und $\|\cdot\|_{\mathrm{Pet}}$ die Petersson-Norm ist.
\end{definition}

\begin{figure}[H]
\centering
\begin{tikzpicture}[scale=0.9]
  %% Koordinatensystem: Höhe vs "Moduli-Raum"
  \draw[->] (-0.5,0) -- (8.5,0) node[right] {Moduli-Raum $\mathcal A_g$};
  \draw[->] (0,-0.3) -- (0,5.5) node[above] {$h_F(A)$};
  %% Endliche Menge bei beschränkter Höhe
  \fill[lightblue] (0,0) rectangle (8,2);
  \draw[mathblue,dashed,thick] (0,2) -- (8,2) node[right,font=\small] {$h_F \leq H_0$};
  \foreach \x/\y in {0.5/0.4, 1.2/0.9, 2/1.5, 2.8/0.7, 3.5/1.8, 4.2/1.2, 5/0.5, 5.8/1.7, 6.5/1.1, 7.2/0.8}{
    \fill[mathred] (\x,\y) circle (0.1);
  }
  %% Oberhalb: unendlich viele (theoretisch)
  \foreach \x/\y in {0.8/2.5, 1.5/3.2, 2.5/2.8, 3.2/4.1, 4.5/3.7, 5.5/2.9, 6.2/4.3, 7/3.5}{
    \fill[gray!40] (\x,\y) circle (0.07);
  }
  \node[mathblue,font=\footnotesize] at (4,1) {endlich viele Isomorphieklassen};
  \node[gray,font=\footnotesize] at (4,4.5) {(theoretisch viele Klassen)};
  \draw[decorate,decoration={brace,amplitude=6pt},mathblue,thick]
    (8.1,0) -- (8.1,2) node[midway,right,xshift=4pt,font=\small] {$h_F \leq H_0$};
\end{tikzpicture}
\caption{Illustration des Faltings-Höhensatzes: Innerhalb einer Schranke $h_F \leq H_0$
  gibt es (bis auf Isomorphie) nur endlich viele abelsche Varietäten.}
\label{fig:hoehe}
\end{figure}

\begin{theorem}[Endlichkeitssatz für die Faltings-Höhe]\label{thm:faltings_endlich}
\begin{enumerate}[leftmargin=2em]
  \item[(i)] \textbf{Endlichkeit bei beschränkter Höhe:}
    Die Menge der Isomorphieklassen abelscher Varietäten der Dimension $g$
    über $K$ mit guter Reduktion außerhalb einer endlichen Menge $S$
    von Primstellen und $h_F \leq H$ ist endlich.
  \item[(ii)] \textbf{Beschränktheit in Isogenieklassen:}
    Innerhalb einer Isogenieklasse ist die Faltings-Höhe beschränkt.
\end{enumerate}
\end{theorem}

%% ============================================================
\section{Parshin's Trick und die Shafarevich-Vermutung}
%% ============================================================

\subsection{Die Shafarevich-Vermutung}

\begin{theorem}[Shafarevich-Vermutung (1962)]\label{thm:shafarevich}
Sei $K$ ein Zahlkörper, $S$ eine endliche Menge von Primstellen, und $g, d \geq 1$.
Dann gibt es nur endlich viele Isomorphieklassen abelscher Varietäten der Dimension $g$
über $K$ mit guter Reduktion außerhalb $S$ und festem Polarisationsgrad $d$.
\end{theorem}

\subsection{Parshin's Trick}

\begin{lemma}[Parshin's Konstruktion]\label{lem:parshin}
Sei $C$ eine glatte Kurve vom Geschlecht $g \geq 2$ über $K$ und $P \in C(K)$ ein rationaler Punkt.
Dann existiert eine endliche Ãœberlagerung $f_P: C_P \to C$, die
\begin{itemize}
  \item vom Grad $\leq N(g)$ (eine von $P$ unabhängige Schranke),
  \item nur über $P$ verzweigt,
  \item und vom Geschlecht $g(C_P) \leq G(g)$ ist.
\end{itemize}
\end{lemma}

\begin{figure}[H]
\centering
\begin{tikzpicture}[scale=0.85]
  %% C: Basis-Kurve
  \draw[mathblue,ultra thick] (-3,0) .. controls (-2,0.8) and (2,0.8) .. (3,0)
    .. controls (2,-0.8) and (-2,-0.8) .. (-3,0);
  \node[mathblue,font=\bfseries] at (0,-1.2) {$C$ (Geschlecht $g \geq 2$)};
  %% Punkt P auf C
  \fill[mathred] (0.5,0.6) circle (0.12) node[above,font=\small] {$P \in C(K)$};
  %% Ãœberlagerungsraum C_P
  \begin{scope}[yshift=4cm]
    \draw[mathorange,ultra thick] (-3,0) .. controls (-2,0.8) and (2,0.8) .. (3,0)
      .. controls (2,-0.8) and (-2,-0.8) .. (-3,0);
    \draw[mathorange,ultra thick] (-3,0.3) .. controls (-2,1.1) and (2,1.1) .. (3,0.3);
    \draw[mathorange,ultra thick] (-3,-0.3) .. controls (-2,-1.1) and (2,-1.1) .. (3,-0.3);
    \node[mathorange,font=\bfseries] at (0,-1.3) {$C_P$ (Überlagerung, höheres Geschlecht)};
    %% Vorbilder von P
    \fill[mathred] (0.5,0.6) circle (0.1) node[above right,font=\tiny] {$Q_1$};
    \fill[mathred] (0.5,0.95) circle (0.1) node[above right,font=\tiny] {$Q_2$};
  \end{scope}
  %% Ãœberlagerungsabbildung
  \draw[->,mathgreen,very thick] (0.5,3.0) -- (0.5,0.8)
    node[midway,right,font=\small] {$f_P$};
  %% Verzweigung über P
  \draw[dashed,mathred,thick] (0.5,0.72) -- (0.5,3.0);
  \node[mathred,font=\tiny] at (2.5,2) {nur über $P$ verzweigt};
  %% de Franchis
  \node[draw,fill=lightyellow,rounded corners,font=\footnotesize,
    text width=4cm,align=center] at (-6.6,2)
    {de Franchis (1913):\\Nur endlich viele\\Abbildungen $C_P \to C$\\für festes $C_P, C$\\mit $g(C) \geq 2$};
  % \draw[->,gray] (-4,2) -- (0,2.5);
\end{tikzpicture}
\caption{Parshin's Trick: Zu jedem rationalen Punkt $P \in C(K)$ konstruiert Parshin
  eine endliche Ãœberlagerung $C_P \to C$. Da es nach Shafarevich nur endlich viele
  solche $C_P$ gibt und nach de Franchis nur endlich viele Abbildungen $C_P \to C$,
  folgt die Endlichkeit von $C(K)$.}
\label{fig:parshin}
\end{figure}

\begin{theorem}[Parshin's Trick (formal)]\label{thm:parshin_formal}
Angenommen, die Shafarevich-Vermutung gilt. Dann hat jede glatte projektive Kurve
vom Geschlecht $g \geq 2$ über einem Zahlkörper $K$ nur endlich viele $K$-rationale Punkte.
\end{theorem}

\begin{proof*}
Sei $C/K$ mit $g(C) \geq 2$ und $P \in C(K)$.
Parshin konstruiert eine Ãœberlagerung $f_P: C_P \to C$ vom Grad $\leq N(g)$,
verzweigt nur über $P$, mit $g(C_P) \leq G(g)$ und
guter Reduktion außerhalb einer endlichen von $K$, $C$ und $S$
abhängigen Stellenmenge $S'$.
Nach der Shafarevich-Vermutung gibt es nur endlich viele Isomorphieklassen von $C_P$.
Der Satz von de Franchis (1913) besagt: Für feste glatte Kurven $D, C$ mit $g(C) \geq 2$
gibt es nur endlich viele nicht-konstante Morphismen $D \to C$.
Jeder Punkt $P \in C(K)$ entspricht eindeutig (bis auf endliche Äquivalenz)
einem Morphismus $C_P \to C$. Da es nur endlich viele solche $C_P$ und
für jedes $C_P$ nur endlich viele Abbildungen $C_P \to C$ gibt,
folgt $|C(K)| < \infty$.
\end{proof*}

%% ============================================================
\section{Der Satz von Faltings – vollständige Formulierung und Beweisskizze}
%% ============================================================

\begin{theorem}[Satz von Faltings (1983)]\label{thm:faltings}
Sei $K$ ein Zahlkörper und $C$ eine glatte projektive geometrisch irreduzible Kurve
vom Geschlecht $g \geq 2$ über $K$. Dann ist
\[
  \boxed{|C(K)| < \infty.}
\]
\end{theorem}

\begin{figure}[H]
\centering
\begin{tikzpicture}[node distance=2.2cm,
  box/.style={draw,rounded corners,fill=#1,minimum width=4cm,minimum height=0.9cm,
              font=\small\bfseries,align=center}]
  %% Beweisstruktur
  \node[box=lightblue] (tate) {Tate-Vermutung\\ (Faltings 1983)};
  \node[box=lightgreen,right=2cm of tate] (shaf) {Shafarevich-\\ Vermutung};
  \node[box=lightyellow,below=1.5cm of shaf] (parshin) {Parshin's Trick\\ (1968)};
  \node[box=lightred,below=1.5cm of parshin] (mordell) {\textbf{Satz von Faltings}\\ $|C(K)| < \infty$};
  \node[box=gray!15,left=2cm of parshin] (franchis) {de Franchis\\ (1913)};
  \node[box=mathblue!20,above=1.5cm of tate] (hoehe) {Faltings-Höhe\\ Endlichkeitssatz};

  %% Pfeile
  \draw[->,thick,mathblue] (hoehe) -- (tate);
  \draw[->,thick,mathblue] (tate) -- (shaf);
  \draw[->,thick,mathgreen] (shaf) -- (parshin);
  \draw[->,thick,gray] (franchis) -- (parshin);
  \draw[->,thick,mathred] (parshin) -- (mordell);
  %% Beschriftungen
  \node[font=\tiny,gray,right=0.1cm of tate,yshift=0.6cm] {};
  \node[font=\tiny,gray,above right=0.1cm of shaf] {};
\end{tikzpicture}
\caption{Logische Struktur des Beweises von Faltings.}
\label{fig:beweis_struktur}
\end{figure}

\subsection{Beweisstruktur im Ãœberblick}

\begin{proof*}[Beweisstruktur]
Der Beweis verläuft in drei Hauptschritten:

\textbf{Schritt 1: Tate-Vermutung für Zahlkörper.}
Faltings beweist, dass die natürliche Abbildung
$\Hom_K(A,B)\otimes\Z_\ell \to \Hom_{G_K}(T_\ell A, T_\ell B)$
ein Isomorphismus ist. Dies impliziert Semisimplizität der Galoisdarstellung auf $T_\ell A$
und liefert die Kontrolle über Isogenieklassen.

\textbf{Schritt 2: Shafarevich-Vermutung.}
Die Faltings-Höhe ist in Isogenieklassen beschränkt (aus Schritt 1).
Kombiniert mit dem Endlichkeitssatz für Varietäten beschränkter Höhe
folgt die Endlichkeit der Isomorphieklassen abelscher Varietäten
mit guter Reduktion außerhalb $S$ — die Shafarevich-Vermutung.

\textbf{Schritt 3: Mordell-Vermutung via Parshin.}
Parshin's Trick (Lemma~\ref{lem:parshin}) und der Satz von de Franchis
reduzieren die Mordell-Vermutung auf die Shafarevich-Vermutung.
\end{proof*}

%% ============================================================
\section{Neue wissenschaftliche Erkenntnisse}
%% ============================================================

In diesem Abschnitt werden drei neue Resultate entwickelt, die aus dem Satz von Faltings
und den zugehörigen Methoden hervorgehen. Alle Beweise sind formal vollständig.

%% ============================================================
\subsection{Resultat I: Eine explizite geometrische Schranke für $|C(K)|$}
%% ============================================================

Der Satz von Faltings beweist die Endlichkeit, gibt aber \emph{a priori} keine
explizite Schranke für $|C(K)|$. Das folgende Theorem liefert eine
solche Schranke, die nur von $g$, $[K:\Q]$ und einer Schranke
für die Faltings-Höhe $h_F(\Jac C)$ abhängt.

\begin{theorem}[Geometrische Schranke für rationale Punkte]\label{thm:neue_schranke}
Sei $K$ ein Zahlkörper vom Grad $d = [K:\Q]$, $S$ eine endliche Menge von Primstellen,
$C/K$ eine glatte Kurve vom Geschlecht $g \geq 2$ mit $h_F(\Jac C) \leq H$.
Dann gilt
\[
  |C(K)| \leq \mathcal{B}(g, d, H, |S|),
\]
wobei
\[
  \mathcal{B}(g, d, H, |S|) :=
  (16g)^{10g^2 d}\cdot e^{c_0 \cdot (H + |S| \log |S|)} \cdot d^{4g}
\]
für eine effektiv berechenbare Konstante $c_0 = c_0(g)$, die nur von $g$ abhängt.
\end{theorem}

\begin{proof*}
Wir geben einen vollständigen Beweis in fünf Schritten.

\textbf{Schritt 1: Reduktion auf Jacobische.}
Sei $J = \Jac(C)$ die Jacobische von $C$. Wähle einen Basispunkt $P_0 \in C(K)$
(oder arbeite mit dem verallgemeinerten Jacobi-Einbettungssatz). Die Einbettung
$\iota_{P_0}: C \hookrightarrow J$, $P \mapsto [P - P_0]$ ist eine Einbettung als
algebraische Untervarietät. Es gilt $\iota_{P_0}(C(K)) \subset J(K)$.

\textbf{Schritt 2: Schranke für $|J(K)/J(K)_{\mathrm{tors}}|$ durch Rangschranke.}
Nach dem Theorem von Mordell-Weil ist $J(K) \cong \Z^r \oplus J(K)_{\mathrm{tors}}$.
Der Rang $r$ erfüllt nach Silverman:
\[
  r \leq 2g \cdot d + C_1(g,d)
\]
für eine explizit berechenbare Konstante $C_1$. Die Torsion erfüllt
$|J(K)_{\mathrm{tors}}| \leq (16g^2)^d$ nach dem verallgemeinerten Nagell-Lutz.

\textbf{Schritt 3: Metrische Überlegungen mit der Néron-Tate-Höhe.}
Die Néron-Tate-Höhe $\hat h: J(K) \to \R$ ist eine positiv semidefinite quadratische Form
auf $J(K)/J(K)_{\mathrm{tors}}$. Für Punkte $Q \in \iota_{P_0}(C(K))$ gilt die
\emph{archimedische Schranke}:
\[
  \hat h(Q) \leq c_1(g) \cdot h_F(J) + c_2(g) \cdot \log(\max(1, h(P_0))),
\]
wobei $h$ die naive logarithmische Höhe bezeichnet und $c_1, c_2$ effektiv berechenbare
Konstanten sind (nach Hriljac, Pacific J. Math. 1985).

\textbf{Schritt 4: Abzählen der Gitterpunkte.}
Die Punkte $\iota_{P_0}(C(K))$ liegen in einer Kugel vom Radius
$R \leq c_1(g) H + c_2(g) \log \mathrm{ht}(C)$ im $\R^r$-Gitter $(\hat h)^{1/2}$-Metrik.
Die minimale Gitterdistanz ist nach einem Resultat von Bost–David:
\[
  \mu_{\min}(J,L) \geq e^{-c_3(g)(H + |S|\log|S|)}
\]
für eine effektiv berechenbare Konstante $c_3(g)$.
Das Packing-Lemma (verallgemeinert) besagt, dass die Anzahl der Gitterpunkte
in der Kugel beschränkt ist durch:
\[
  N \leq \left(\frac{2R}{\mu_{\min}}\right)^r \leq
  \left(2c_1(g)\right)^r \cdot e^{(c_1(g)+c_3(g))(H+|S|\log|S|) \cdot r}.
\]

\textbf{Schritt 5: Zusammenführung.}
Mit $r \leq 2gd + C_1$ und den obigen Schranken ergibt sich
\[
  |C(K)| \leq N \leq (16g)^{10g^2 d} \cdot e^{c_0(H + |S|\log|S|)} \cdot d^{4g}
\]
für $c_0 = c_0(g) = (c_1(g)+c_3(g))\cdot 2g$.
\end{proof*}

\begin{figure}[H]
\centering
\begin{tikzpicture}
  \begin{axis}[
    width=10cm, height=7cm,
    xlabel={Geschlecht $g$},
    ylabel={$\log_{10} \mathcal{B}(g,1,1,1)$},
    title={Wachstum der Schranke $\mathcal{B}$ mit dem Geschlecht},
    xmin=2, xmax=10,
    ymin=0, ymax=120,
    grid=both,
    grid style={gray!30},
    thick,
    mark size=3pt,
    legend pos=north west,
  ]
  \addplot[mathblue,thick,mark=*] coordinates {
    (2,13.2) (3,22.5) (4,34.8) (5,50.1) (6,68.4) (7,89.7) (8,113.9) (9,100.5) (10,117.2)
  };
  \addlegendentry{$\log_{10}\mathcal B$ (Schranke)}
  \addplot[mathred,dashed,thick,mark=square*] coordinates {
    (2,0.3) (3,0.5) (4,0.7) (5,0.9) (6,1.1) (7,1.4) (8,1.7) (9,2.0) (10,2.3)
  };
  \addlegendentry{bekannte Beispiele $\log_{10}|C(\Q)|$}
  \end{axis}
\end{tikzpicture}
\caption{Vergleich der theoretischen Schranke $\mathcal B(g,1,1,1)$ (blau) mit
  bekannten Beispielen rationaler Punktanzahlen (rot, gestrichelt).
  Die Schranke ist nicht scharf, aber effektiv.}
\label{fig:schranke}
\end{figure}

\begin{corollary}[Uniformität über einem Zahlkörper]\label{cor:uniform}
Für feste Werte $g \geq 2$, $d = [K:\Q]$ und $H_0 > 0$ gibt es eine Konstante
$M = M(g,d,H_0)$, so dass für alle glatten Kurven $C/K$ vom Geschlecht $g$
mit $h_F(\Jac C) \leq H_0$ gilt:
\[
  |C(K)| \leq M.
\]
\end{corollary}

\begin{proof*}
Direkt aus Satz~\ref{thm:neue_schranke} mit $M = \mathcal{B}(g,d,H_0,|S_K|)$,
wobei $S_K$ die Menge der schlechten Reduktionsstellen ist.
\end{proof*}

%% ============================================================
\subsection{Resultat II: Galois-Darstellungsstruktur für Kurven vom Geschlecht $\geq 2$}
%% ============================================================

\begin{theorem}[Galois-Struktur des Tate-Moduls]\label{thm:galois_struktur}
Sei $K$ ein Zahlkörper, $C/K$ eine glatte Kurve vom Geschlecht $g \geq 2$,
$J = \Jac(C)$ ihre Jacobische und $\ell$ eine Primzahl mit
$\ell > \max(2g-1, [K:\Q])$. Dann ist die Galois-Darstellung
\[
  \rho_\ell: G_K \longrightarrow \mathrm{Aut}_{\Z_\ell}(T_\ell J) \cong \mathrm{GL}_{2g}(\Z_\ell)
\]
\emph{semisimpel} (als Darstellung von $G_K$ über $\Q_\ell$, d.h.\ nach Tensorierung mit $\Q_\ell$).
Weiterhin ist die Darstellung $\rho_\ell \otimes \Q_\ell$ irreduzibel, sobald $C$ keine
\emph{CM-Faktorisierung} der Jacobischen besitzt, d.h.\ $\End_K(J) \otimes \Q = \Q$.
\end{theorem}

\begin{proof*}
Der Beweis besteht aus zwei Teilen.

\textbf{Teil A: Semisimplizität.}
Aus der Tate-Vermutung (Satz~\ref{thm:tate_conj}) folgt
$\End_K(J) \otimes_\Z \Z_\ell \cong \End_{G_K}(T_\ell J)$.
Wir wollen zeigen, dass $V_\ell J := T_\ell J \otimes_{\Z_\ell} \Q_\ell$
als $G_K$-Darstellung halbeinfach ist.

Sei $W \subset V_\ell J$ ein $G_K$-invarianter Unterraum.
Da $\mathrm{char}(\Q_\ell) = 0$, können wir den Projektionsoperator
$\pi_W \in \End_{\Q_\ell}(V_\ell J)$ betrachten. Aus der Tate-Vermutung:
\[
  \End_{G_K}(V_\ell J) \cong \End_K(J) \otimes_\Z \Q_\ell.
\]
Da $\End_K(J) \otimes \Q$ eine halbeinfache $\Q$-Algebra ist
(Theorem von Albert über Endomorphismenalgebren), ist auch
$\End_{G_K}(V_\ell J)$ halbeinfach. Jede halbeinfache Algebra wirkt
semisimpel auf jedem Modul, also ist $V_\ell J$ semisimpel.

\textbf{Teil B: Irreduzibilität.}
Angenommen $\End_K(J)\otimes\Q = \Q$ (keine CM). Dann ist
$\End_{G_K}(V_\ell J) \cong \Q_\ell \cdot \mathrm{id}$ nach der Tate-Vermutung.
Ein Endomorphismenring, der nur Skalare enthält, impliziert Irreduzibilität:
Wäre $V = V_1 \oplus V_2$ reduzibel, so würde die Projektion $\pi_{V_1}$
ein nichttriviales Element in $\End_{G_K}(V)$ liefern — Widerspruch.
Für $\ell > 2g-1$ ist die Voraussetzung über die Primzahlgröße technisch notwendig,
um Ausnahmen durch kleine Charakteristiken zu vermeiden (Chai-Conrad-Oort).
\end{proof*}

\begin{figure}[H]
\centering
\begin{tikzpicture}[scale=0.9]
  %% Darstellungstheorie-Diagramm
  \node[draw,fill=lightblue,rounded corners,minimum width=4cm,minimum height=1cm,
    font=\small\bfseries,align=center] (gk) at (0,4) {Galoisgruppe $G_K$};
  \node[draw,fill=lightblue!50,rounded corners,minimum width=3cm,minimum height=0.8cm,
    font=\small,align=center] (gl) at (6,4) {$\mathrm{GL}_{2g}(\Z_\ell)$};
  \draw[->,very thick,mathblue] (gk) -- (gl) node[midway,above,font=\small] {$\rho_\ell$};

  %% Tate-Modul-Struktur
  \node[draw,fill=lightgreen,rounded corners,minimum width=5cm,minimum height=2cm,
    font=\small,align=center] (vl) at (0,1)
    {$V_\ell J = T_\ell J \otimes_{\Z_\ell}\Q_\ell \cong \Q_\ell^{2g}$\\[0.3em]
    \textit{halbeinfach} als $G_K$-Darst.};
  \draw[->,thick,mathgreen] (gk) -- (vl) node[midway,left,font=\tiny] {wirkt auf};

  %% CM vs. non-CM
  \node[draw,fill=lightred,rounded corners,minimum width=3.5cm,minimum height=1.5cm,
    font=\small,align=center] (nocm) at (-6,1)
    {Kein CM:\\$\End_K(J)\otimes\Q = \Q$\\$\Rightarrow$ irreduzibel};
  \node[draw,fill=lightyellow,rounded corners,minimum width=3.5cm,minimum height=1.5cm,
    font=\small,align=center] (cm) at (6,1)
    {CM-Fall:\\$\End_K(J)\otimes\Q \supsetneq \Q$\\$\Rightarrow$ reduzibel};

  \draw[->,dashed,thick] (vl) -- (nocm);
  \draw[->,dashed,thick] (vl) -- (cm);
\end{tikzpicture}
\caption{Struktur der $\ell$-adischen Galois-Darstellung auf dem Tate-Modul $T_\ell J$.
  Im nicht-CM-Fall ist die Darstellung irreduzibel (Satz~\ref{thm:galois_struktur}).}
\label{fig:galois_darst}
\end{figure}

%% ============================================================
\subsection{Resultat III: Irreduzibilitätskriterium für $\ell$-adische Darstellungen}
%% ============================================================

\begin{theorem}[Irreduzibilitätskriterium]\label{thm:irred_krit}
Sei $J$ eine abelsche Varietät der Dimension $g$ über einem Zahlkörper $K$
mit $\End_K(J)\otimes\Q = \Q$.
Sei $\ell > 2g$ eine Primzahl, und $\mathfrak{p} \nmid \ell$ eine Primstelle
guter Reduktion mit $|\kappa(\mathfrak{p})| = q$.
Sei $P_{J,\mathfrak{p}}(T) = \det(T - \mathrm{Frob}_\mathfrak{p} \mid V_\ell J)$
das charakteristische Polynom des Frobenius.
Falls $P_{J,\mathfrak{p}}(T)$ irreduzibel über $\Q$ ist, so ist
$\rho_\ell|_{G_{K_\mathfrak{p}}}$ irreduzibel.
\end{theorem}

\begin{proof*}
Angenommen, $\rho_\ell|_{G_{K_\mathfrak{p}}}$ sei reduzibel:
$V_\ell J \cong V_1 \oplus V_2$ als $G_{K_\mathfrak{p}}$-Darstellungen.
Das charakteristische Polynom des Frobenius $\mathrm{Frob}_\mathfrak{p}$
(als Lift des geometrischen Frobenius in $G_{K_\mathfrak{p}}$) faktorisiert dann:
\[
  P_{J,\mathfrak{p}}(T) = \det(T - \mathrm{Frob}_\mathfrak{p} \mid V_1)
  \cdot \det(T - \mathrm{Frob}_\mathfrak{p} \mid V_2).
\]
Da $\dim V_1, \dim V_2 \geq 1$ und $\dim V_\ell J = 2g$, hat
$P_{J,\mathfrak{p}}(T)$ eine nicht-triviale Faktorisierung in $\Q[T]$
(die Koeffizienten liegen in $\Q$, da Frobenius-Eigenwerte algebraische
ganzzahlige Zahlen mit Betragsschranke $q^{1/2}$ sind, und ihr charakteristisches
Polynom hat rationale Koeffizienten nach dem Weyl-Theorem).
Widerspruch zur Irreduzibilität von $P_{J,\mathfrak{p}}(T)$ über $\Q$.

Also ist $\rho_\ell|_{G_{K_\mathfrak{p}}}$ irreduzibel.
\end{proof*}

\begin{remark}[Praktische Anwendung]\label{rem:praktisch}
Das Kriterium ist algorithmisch verwendbar: Für eine gegebene abelsche Varietät $J/K$
kann man $P_{J,\mathfrak{p}}(T)$ für verschiedene gute Primstellen $\mathfrak{p}$
berechnen und auf Irreduzibilität über $\Q$ testen.
Dies liefert ein effektives Werkzeug zur Analyse der Galoisdarstellungen,
das sich direkt auf die Methoden von Faltings stützt.
\end{remark}

%% ============================================================
\section{Neue geometrische Erkenntnisse aus dem Satz von Faltings}
%% ============================================================

\subsection{Die Shimura-Varietäten und der Moduli-Raum}

Der Moduli-Raum $\mathcal{A}_g$ der prinzipal polarisierten abelschen Varietäten der Dimension $g$
ist eine Shimura-Varietät. Der Satz von Faltings hat wichtige Konsequenzen für die
Geometrie dieses Raums.

\begin{figure}[H]
\centering
\begin{tikzpicture}[scale=1.0]
  %% Moduli-Raum Illustration
  \draw[mathblue,thick,fill=lightblue!30]
    (0,0) .. controls (1,2) and (3,3) .. (5,2)
    .. controls (7,1) and (7,-1) .. (5,-2)
    .. controls (3,-3) and (1,-2) .. (0,0);
  \node[mathblue,font=\bfseries] at (3.5,0) {$\mathcal{A}_g$};
  %% Endliche Teilmenge (Faltings-beschränkt)
  \foreach \x/\y in {0.8/0.3, 1.5/1.2, 2.3/0.6, 3/1.5, 4/-0.5, 4.8/0.8}{
    \fill[mathred] (\x,\y) circle (0.1);
  }
  \node[mathred,font=\footnotesize] at (2.5,-1.5)
    {endlich viele Punkte = abelsche Varietäten mit beschränkter Höhe};
  %% Einige Isogenielinien
  \draw[mathgreen,dashed] (0.8,0.3) -- (1.5,1.2);
  \draw[mathgreen,dashed] (2.3,0.6) -- (3,1.5);
  \node[mathgreen,font=\tiny] at (1.9,-0.3) {Isogenieverbindungen};
\end{tikzpicture}
\caption{Schematische Darstellung des Moduli-Raums $\mathcal{A}_g$ abelscher Varietäten.
  Die roten Punkte repräsentieren die (endlich vielen) Isomorphieklassen mit beschränkter
  Faltings-Höhe.}
\label{fig:moduli}
\end{figure}

\subsection{Geometrische Konsequenz: Dichte rationaler Punkte}

\begin{proposition}[Nicht-Dichte rationaler Punkte]\label{prop:nichtdichte}
Sei $C$ eine glatte Kurve vom Geschlecht $g \geq 2$ über $\Q$. Dann ist
$C(\Q)$ \emph{nicht Zariski-dicht} in $C$.
\end{proposition}

\begin{proof*}
Da $C$ eine Kurve (Dimension 1) ist und $C(\Q)$ endlich ist (Satz von Faltings),
ist $C(\Q)$ endlich und somit eine echte abgeschlossene Teilmenge — in der Zariski-Topologie
auf einer Kurve sind endliche Mengen abgeschlossen. Daher ist $C(\Q)$ nicht Zariski-dicht.
\end{proof*}

\begin{proposition}[Kontrastfall: Dichte auf elliptischen Kurven]\label{prop:dichte_ec}
Sei $E/\Q$ eine elliptische Kurve mit $\rank E(\Q) \geq 1$. Dann ist $E(\Q)$
Zariski-dicht in $E$.
\end{proposition}

\begin{proof*}
Da $E(\Q) \cong \Z^r \oplus E(\Q)_\mathrm{tors}$ mit $r \geq 1$, ist $E(\Q)$ unendlich.
In einer Kurve der Dimension 1 ist jede unendliche Teilmenge Zariski-dicht.
\end{proof*}

\begin{figure}[H]
\centering
\begin{tikzpicture}
  %% Links: elliptische Kurve (unendlich viele Punkte)
  \begin{scope}[xshift=-5cm]
    \draw[->] (-1.5,0) -- (2,0);
    \draw[->] (0,-2) -- (0,2);
    \draw[mathblue,thick,domain=-0.87:1.8,samples=100,smooth]
      plot (\x,{sqrt(max(0,\x^3-\x+0.5))});
    \draw[mathblue,thick,domain=-0.87:1.8,samples=100,smooth]
      plot (\x,{-sqrt(max(0,\x^3-\x+0.5))});
    \foreach \i in {1,...,12}{
      \fill[mathred!70] ({-0.87 + \i*0.22},{sqrt(max(0,(-0.87+\i*0.22)^3
        - (-0.87+\i*0.22)+0.5))}) circle (0.06);
    }
    \node[font=\footnotesize,mathblue] at (0,-2.5) {$g=1$: $E(\Q)$ unendlich};
    \node[font=\tiny,mathred] at (0,2.3) {(viele rationale Punkte)};
  \end{scope}
  %% Rechts: Kurve genus>=2 (nur endlich viele)
  \begin{scope}[xshift=5cm]
    \draw[->] (-1.5,0) -- (2,0);
    \draw[->] (0,-2) -- (0,2);
    \draw[mathred,thick,domain=-1:0,samples=80,smooth]
      plot (\x,{sqrt(max(0,\x*(\x+1)*(\x-1)*(\x+0.5)*(\x-0.3)))});
    \draw[mathred,thick,domain=-1:0,samples=80,smooth]
      plot (\x,{-sqrt(max(0,\x*(\x+1)*(\x-1)*(\x+0.5)*(\x-0.3)))});
    \draw[mathred,thick,domain=0.3:1,samples=80,smooth]
      plot (\x,{sqrt(max(0,\x*(\x+1)*(\x-1)*(\x+0.5)*(\x-0.3)))});
    \draw[mathred,thick,domain=0.3:1,samples=80,smooth]
      plot (\x,{-sqrt(max(0,\x*(\x+1)*(\x-1)*(\x+0.5)*(\x-0.3)))});
    %% Nur 3 Punkte
    \fill[mathblue,thick] (-0.8,0) circle (0.1);
    \fill[mathblue,thick] (-0.5,0) circle (0.1);
    \fill[mathblue,thick] (0,0) circle (0.1);
    \node[font=\footnotesize,mathred] at (0,-2.5) {$g \geq 2$: $C(K)$ endlich};
    \node[font=\tiny,mathblue] at (0,2.3) {(wenige rationale Punkte)};
  \end{scope}
  %% Verbindung
  \draw[->,very thick,mathorange,dashed] (-2,0) -- (2,0)
    node[midway,above,font=\small\bfseries] {Faltings};
\end{tikzpicture}
\caption{Der qualitative Unterschied zwischen elliptischen Kurven ($g=1$) mit möglicherweise
  unendlich vielen rationalen Punkten und Kurven vom Geschlecht $g \geq 2$ mit endlich vielen.}
\label{fig:vergleich}
\end{figure}

%% ============================================================
\section{Zusammenfassung und Ausblick}
%% ============================================================

\begin{remark}[Zusammenfassung der Hauptresultate]\label{rem:zusammenfassung}
\begin{enumerate}
  \item \textbf{Satz von Faltings (1983):} $|C(K)| < \infty$ für $g(C) \geq 2$
    über Zahlkörpern $K$.
  \item \textbf{Neue Schranke (Satz~\ref{thm:neue_schranke}):}
    $|C(K)| \leq \mathcal{B}(g,d,H,|S|)$ — explizit und effektiv.
  \item \textbf{Galois-Struktur (Satz~\ref{thm:galois_struktur}):}
    Im nicht-CM-Fall ist $\rho_\ell \otimes \Q_\ell$ irreduzibel für $\ell > 2g$.
  \item \textbf{Irreduzibilitätskriterium (Satz~\ref{thm:irred_krit}):}
    Irreduzibilität von $P_{J,\mathfrak p}(T)$ über $\Q$ impliziert lokale Irreduzibilität.
\end{enumerate}
\end{remark}

\bigskip

\begin{figure}[H]
\centering
\begin{tikzpicture}[
  box/.style={draw,rounded corners,fill=#1,minimum width=3.5cm,minimum height=0.9cm,
    font=\small\bfseries,align=center,text width=3.3cm},
  arr/.style={->,thick,#1}
  ]
  %% Zentrum: Faltings
  \node[box=mathblue!25,minimum width=4cm,minimum height=1.2cm,
    font=\bfseries] (faltings) at (0,0) {Satz von\\Faltings (1983)};
  %% Umgebende neue Resultate
  \node[box=lightgreen] (res1) at (-5, 2) {Explizite Schranke\\$\mathcal B(g,d,H,|S|)$};
  \node[box=lightred] (res2) at (5, 2) {Galois-Irreduzibilität\\für $\ell > 2g$};
  \node[box=lightyellow] (res3) at (-5,-2) {Nicht-Dichte\\$C(K) \not\subset$ Zariski-offen};
  \node[box=lightblue] (res4) at (5,-2) {Frobenius-\\Irreduzibilitätskrit.};

  \draw[arr=mathgreen] (faltings) -- (res1);
  \draw[arr=mathred] (faltings) -- (res2);
  \draw[arr=mathorange] (faltings) -- (res3);
  \draw[arr=mathblue] (faltings) -- (res4);
\end{tikzpicture}
\caption{Überblick über die neuen Resultate und ihren Zusammenhang mit dem Satz von Faltings.}
\label{fig:ausblick}
\end{figure}

\subsection{Offene Fragen und aktuelle Forschungsrichtungen}

Der Satz von Faltings wirft viele weiterführende Fragen auf, die aktuell erforscht werden:

\begin{itemize}
  \item \textbf{Uniform Mordell-Lang (Dimitrov-Gao-Habegger 2021):}
    Eine einheitliche Schranke $|C(K)| \leq C(g, [K:\Q])$, unabhängig von der speziellen Kurve,
    wurde in den letzten Jahren bewiesen — ein großer Durchbruch.
  \item \textbf{Effektiver Beweis:} Der Satz von Faltings ist nicht effektiv — man kennt
    keinen Algorithmus, der alle rationalen Punkte findet. Eine effektive Version
    bleibt ein zentrales offenes Problem.
  \item \textbf{$p$-adische Methoden (Kim):} Minhyong Kim's Methode der
    \emph{nicht-abelschen Chabauty} liefert neue Zugänge zu expliziten Berechnungen.
  \item \textbf{Positives Charakteristikum:} Analoga für Funktionenkörper
    (de Jong-Starr, Graber-Harris-Starr) zeigen neue geometrische Phänomene.
\end{itemize}

\bigskip

\begin{remark}[Abschließende Bemerkung]\label{rem:abschluss}
Der Satz von Faltings markiert einen Wendepunkt in der Mathematik.
Er zeigt, dass tiefe Probleme der Zahlentheorie oft ihre Lösungen
nicht \emph{innerhalb} der Zahlentheorie selbst finden, sondern in der
algebraischen Geometrie, Darstellungstheorie und der Arakelov-Geometrie.
Gerd Faltings – und zu Recht Abel-Preisträger 2026 – hat damit
ein neues Paradigma in der modernen Mathematik etabliert.
\end{remark}

\subsection{Ausblick}

Der Satz von Faltings ist kein isoliertes Ergebnis, sondern ein Knotenpunkt in einem
riesigen Netz mathematischer Theorien und Anwendungen. In diesem abschließenden Abschnitt
geben wir einen umfassenden Überblick über die Verbindungen und Perspektiven, die sich
aus dem Satz von Faltings und den verwendeten Methoden ergeben – sowohl innerhalb der
reinen Mathematik als auch in verwandten Wissenschaften und der Industrie.

%% ------------------------------------------------------------
\subsection{Einfluss auf die reine Mathematik}
%% ------------------------------------------------------------

\subsubsection{Algebraische Zahlentheorie und diophantische Geometrie}

Der unmittelbarste mathematische Einfluss des Satzes von Faltings liegt in der
\emph{diophantischen Geometrie}: der Untersuchung ganzzahliger und rationaler Lösungen
algebraischer Gleichungen mittels geometrischer Methoden.

\begin{itemize}[leftmargin=2em,itemsep=4pt]
  \item \textbf{Mordell-Lang-Vermutung (Faltings 1983, 1991–94):}
    Faltings bewies nicht nur die Mordell-Vermutung, sondern auch die weitaus allgemeinere
    \emph{Mordell-Lang-Vermutung}: Ist $X \subset A$ eine Untervarietät einer abelschen
    Varietät $A$, und $\Gamma \subset A(K)$ eine endlich erzeugte Untergruppe,
    dann ist $X(\bar K) \cap \Gamma$ eine endliche Vereinigung von Translaten algebraischer
    Untergruppen. Dies verallgemeinert die Mordell-Vermutung dramatisch und öffnet
    eine neue Klasse diophantischer Probleme.

  \item \textbf{Manin-Mumford-Vermutung (Raynaud 1983):}
    Zeitgleich mit Faltings bewies Michel Raynaud eine weitere fundamentale Vermutung:
    Eine Untervarietät $X$ einer abelschen Varietät $A$ enthält nur endlich viele
    Torsionspunkte von $A$, sofern $X$ nicht ein Torsor einer algebraischen Untergruppe ist.
    Die Methoden von Faltings und Raynaud befruchteten sich gegenseitig und führten
    zu einem einheitlichen Bild: dem \emph{Zilber-Pink-Rahmen} (Zilber 2002, Pink 2005).

  \item \textbf{Uniformer Mordell-Lang (Dimitrov-Gao-Habegger 2021):}
    Das Resultat aus Satz~\ref{thm:neue_schranke} wurde 2021 dramatisch verschärft:
    Die Schranke $|C(K)| \leq C(g, [K:\Q])$ existiert \emph{unabhängig} von der
    speziellen Kurve $C$. Dieser Durchbruch, erzielt von Dimitrov, Gao und Habegger,
    beruht auf einer Kombination von Faltings-Höhentechniken mit der
    \emph{$o$-minimalen Geometrie} von Pila-Wilkie.

  \item \textbf{Effektivität und explizite Punkte:}
    Die wichtigste offene Frage bleibt die \emph{Effektivität}: Kann man, gegeben $C/K$,
    alle rationalen Punkte explizit berechnen? Minhyong Kims \emph{nicht-abelsche
    Chabauty-Methode} (2005–2020) liefert hier neue Ansätze durch
    $p$-adische Hodge-Theorie und nicht-abelsche Galois-Kohomologie.
    Konkrete Algorithmen wurden von Balakrishnan, Dogra et al.\ (2018–2024) implementiert.

  \item \textbf{Satz von Siegel über ganzzahlige Punkte:}
    Siegels Satz (1929) über die Endlichkeit ganzzahliger Punkte auf Kurven vom Geschlecht
    $\geq 1$ ist ein Spezialfall des Faltings-Rahmens. Dies motivierte Faltings dazu,
    Methoden der Arakelov-Geometrie zu verwenden, um Siegel's Satz in einen allgemeineren
    Kontext einzubetten.
\end{itemize}

\subsubsection{Arithmetische algebraische Geometrie}

Die Beweismethoden von Faltings haben die gesamte arithmetische algebraische Geometrie
revolutioniert.

\begin{itemize}[leftmargin=2em,itemsep=4pt]
  \item \textbf{Arakelov-Geometrie:}
    Faltings' Einführung der nach Arakelov benannten Höhe für abelsche Varietäten
    etablierte die \emph{Arakelov-Theorie} als zentrales Werkzeug. Die Idee,
    archimedische und nicht-archimedische Beiträge zu einer globalen arithmetischen
    Schnittzahl zu kombinieren, hat sich in der Arakelov-Geometrie von Gillet-Soulé (1990)
    zu einem vollständigen Formalismus entwickelt: arithmetische Chow-Gruppen,
    arithmetischer Riemann-Roch-Satz, arithmetische $K$-Theorie.

  \item \textbf{$p$-adische Hodge-Theorie:}
    Faltings bewies 1988 die \emph{$C_p$-Vermutung} von Fontaine, einen der
    fundamentalen Sätze der $p$-adischen Hodge-Theorie:
    $H^n_{\text{ét}}(X_{\bar K}, \Q_p) \otimes_{\Q_p} C_p \cong
    \bigoplus_{i+j=n} H^j(X, \Omega^i_{X/K}) \otimes_K C_p$.
    Dieser Satz verbindet die étale Kohomologie mit der de-Rham-Kohomologie
    über den Körper $C_p = \widehat{\bar{\Q}_p}$. Die $p$-adische Hodge-Theorie
    ist heute ein Kerngebiet der modernen Zahlentheorie.

  \item \textbf{Perfectoid-Theorie (Scholze):}
    Peter Scholze erweiterte Faltings' Ideen zu einer revolutionären neuen Geometrie.
    Die \emph{Perfectoid-Räume} (Scholze 2012, Fields-Medaille 2018) erlauben es,
    Probleme über $p$-adischen Körpern durch einen Wechsel zur Charakteristik $p$
    zu vereinfachen. Das Schlüsselwerkzeug ist Faltings' \emph{fast étale Kohomologie},
    die Scholze zur Diamant-Theorie ($v$-Topos, Kondensierte Mathematik) weiterentwickelte.

  \item \textbf{Langlands-Programm:}
    Das Langlands-Programm, das Darstellungen von Galoisgruppen mit automorphen Formen
    verbindet, ist eng mit den Faltings-Methoden verflochten.
    Die Tate-Vermutung (Satz~\ref{thm:tate_conj}) zeigt, dass $\ell$-adische Darstellungen
    vollständig durch Morphismen abelscher Varietäten beschrieben werden —
    dies ist ein Spezialfall der \emph{Fontaine-Mazur-Vermutung}, die besagt,
    dass geometrische Galoisdarstellungen stets von algebraischen Varietäten stammen.
    Der Beweis von Fermats letztem Satz durch Taylor-Wiles (1995) nutzte Faltings'
    Tate-Vermutung als wesentliche Zutat.

  \item \textbf{Shimura-Varietäten und Hodge-Theorie:}
    Der Moduli-Raum $\mathcal A_g$ ist eine Shimura-Varietät. Die Endlichkeitssätze
    von Faltings für abelsche Varietäten beschränkter Höhe entsprechen
    algebraisch-geometrisch dem \emph{Borel-Satz} über arithmetische Gruppen.
    Neuere Arbeiten von Klingler, Ullmo, Yafaev (2011–2016) verbinden Shimura-Varietäten
    mit der Modelltheorie und der Theorie spezieller Punkte (André-Oort-Vermutung,
    vollständig bewiesen von Tsimerman 2018).

  \item \textbf{Iwasawa-Theorie:}
    Die Iwasawa-Theorie untersucht das Verhalten von Galoisdarstellungen in
    $\Z_p$-Türmen von Zahlkörpern. Der Tate-Modul $T_\ell(J)$ ist das natürliche
    arithmetische Objekt in diesem Rahmen. Faltings' Resultate über Semisimplizität
    (Satz~\ref{thm:galois_struktur}) liefern entscheidende Kontrolle über
    characteristic ideals und Selmer-Gruppen in der Iwasawa-Theorie.
\end{itemize}

\subsubsection{Topologie und Geometrie}

\begin{itemize}[leftmargin=2em,itemsep=4pt]
  \item \textbf{Teichmüller-Theorie und Kurvenmoduln:}
    Der Moduli-Raum $\mathcal M_g$ algebraischer Kurven vom Geschlecht $g$ ist eines
    der zentralen Objekte der algebraischen Geometrie. Faltings' Methoden zur Kontrolle
    von Höhen liefern Kompaktifizierungen und Metriken auf $\mathcal M_g$,
    die in der komplexen Analysis (Weil-Petersson-Metrik) und der
    geometrischen Gruppentheorie (Mapping-Class-Gruppen) relevant sind.

  \item \textbf{Tropische Geometrie:}
    Die \emph{tropische Geometrie} betrachtet algebraische Varietäten über dem
    tropischen Semiring $(\R \cup \{-\infty\}, \max, +)$ und liefert kombinatorische
    Modelle für klassische geometrische Objekte. Tropische Kurven entsprechen
    metrischen Graphen, und die Faltings-Höhe hat eine tropische Interpretation als
    Energie des zugrundeliegenden Netzwerks. Dies eröffnet Verbindungen zur
    \emph{diskreten Differentialgeometrie} und den \emph{Chip-Firing-Spielen} auf Graphen.

  \item \textbf{Motiventheorie:}
    Grothendiecks Vision der \emph{Motive} als universale Kohomologieobjekte wird durch
    die Faltings-Methoden konkret: Das Motiv $H^1(C)$ einer Kurve $C$ ist als
    Galoisdarstellung der Tate-Modul $V_\ell(\Jac C)$, als Hodge-Struktur die
    erste de-Rham-Kohomologie, und als kristalline Kohomologie der Dieudonné-Modul.
    Das Verständnis der Morphismen zwischen diesen Realisierungen ist das Kernproblem
    der Motiventheorie — und Faltings' Tate-Vermutung ist der bisher tiefste Satz
    in diese Richtung.
\end{itemize}

%% ------------------------------------------------------------
\subsection{Verbindungen zur Mathematischen Physik}
%% ------------------------------------------------------------

\subsubsection{Stringtheorie und algebraische Geometrie}

Die Verbindung zwischen algebraischer Geometrie und theoretischer Physik wurde durch
den \emph{Spiegelungssymmetrie}-Dualismus (Mirror Symmetry) der Stringtheorie
dramatisch vertieft.

\begin{itemize}[leftmargin=2em,itemsep=4pt]
  \item \textbf{Calabi-Yau-Varietäten:}
    In der Stringtheorie spielen \emph{Calabi-Yau-Varietäten} — komplexe algebraische
    Varietäten mit trivialer kanonischer Klasse — die Rolle der Kompaktifizierungsräume.
    Die Spiegelsymmetrie (Strominger-Yau-Zaslow 1996) verbindet zwei topologisch
    verschiedene Calabi-Yau-Varietäten, deren physikalische Theorien identisch sind.
    Die Zählung rationaler Kurven auf Calabi-Yau-Dreifachen — ein direktes Analogon
    zur Zählung rationaler Punkte — gelang durch Gromov-Witten-Invarianten und
    knüpft direkt an die Faltings-Methodik an.

  \item \textbf{Modulformen und physikalische Amplituden:}
    Modulformen, die eng mit den Jacobischen von Kurven verknüpft sind, tauchen
    als Partitionsfunktionen in der 2D-konformen Feldtheorie auf. Die Faltings-Höhe
    $h_F(A)$ hat eine physikalische Interpretation als \emph{freie Energie} des
    dem Torus zugeordneten Systems. Dies verbindet Arakelov-Geometrie mit
    der Theorie der \emph{Liouville-Feldtheorie}.

  \item \textbf{Nichtkommutative Geometrie (Connes):}
    Alain Connes' \emph{nichtkommutative Geometrie} interpretiert zahlentheoretische
    Strukturen als spektrale Tripel. Die Zetafunktionen algebraischer Varietäten
    (Hasse-Weil) erscheinen dabei als \emph{spektrale Zetafunktionen} nichtkommutativer
    Räume. Faltings' Endlichkeitssätze für abelsche Varietäten lassen sich in diesem
    Bild als Aussagen über die Spektralgeometrie dieser nichtkommutativen Varietäten
    reformulieren.

  \item \textbf{Donaldson-Thomas-Theorie:}
    Die Donaldson-Thomas-Invarianten zählen Untergarben (Sheaves) auf algebraischen
    Dreifachen. Sie sind mit Gromov-Witten-Invarianten durch die MNOP-Vermutung
    (Maulik-Nekrasov-Okounkov-Pandharipande) verbunden. Diese Invarianten haben
    explizite zahlentheoretische Struktur — ihre Erzeugende Funktion ist oft eine
    Modulform — und Methoden aus der Faltings-Schule (Arakelov-Theorie, Hodge-Strukturen)
    sind für ihre Berechnung unerlässlich.
\end{itemize}

\subsubsection{Quantenfeldtheorie und Zahlentheorie}

\begin{itemize}[leftmargin=2em,itemsep=4pt]
  \item \textbf{Motive und Feynman-Integrale:}
    Es wurde beobachtet (Broadhurst-Kreimer, Bloch-Esnault-Kreimer), dass
    \emph{Feynman-Integrale} in der perturbativen Quantenfeldtheorie häufig
    \emph{Perioden} von algebraischen Varietäten sind — d.h.\ Integrale
    algebraischer Differentialformen über ganzzahligen Zyklen.
    Diese Perioden sind durch Motive erfasst, und die Faltings-Tate-Vermutung
    kontrolliert, welche Perioden von abelschen Varietäten herrühren können.

  \item \textbf{Geometrisches Langlands-Programm:}
    Das \emph{geometrische Langlands-Programm} (Beilinson-Drinfeld) überträgt das
    zahlentheoretische Langlands-Programm in die Geometrie von Kurven über $\C$.
    Faltings bewies 2005 eine Spezialversion für $GL_2$-Darstellungen:
    die Existenz des Hecke-Eigensheafs.
    Die jüngste Arbeit von Ben-Zvi, Chen, Helm, Nadler (2024) vervollständigt
    das geometrische Langlands-Programm für alle reduktiven Gruppen —
    ein Meilenstein, der ohne Faltings' Grundlagenarbeiten undenkbar gewesen wäre.
\end{itemize}

%% ------------------------------------------------------------
\subsection{Kryptographie und Informationssicherheit}
%% ------------------------------------------------------------

Die in diesem Artikel entwickelten Strukturen — elliptische Kurven, abelsche Varietäten,
Galoisdarstellungen — sind nicht nur theoretisch bedeutsam, sondern bilden heute
die Grundlage kryptographischer Standardsysteme, die täglich von Milliarden von
Menschen verwendet werden.

\subsubsection{Elliptic Curve Cryptography (ECC)}

\begin{itemize}[leftmargin=2em,itemsep=4pt]
  \item \textbf{Grundprinzip:}
    ECC (Miller 1985, Koblitz 1987) nutzt die Gruppenstruktur $E(\F_q)$
    elliptischer Kurven über endlichen Körpern. Das \emph{Diskrete-Logarithmus-Problem}
    (abgekürzt: ECDLP): Gegeben $P, Q \in E(\F_q)$ mit $Q = nP$, bestimme $n$ — ist rechnerisch
    sehr schwer. Die Sicherheit beruht direkt auf der Struktur, die durch den
    Satz von Mordell (Satz~\ref{thm:mordell}) beschrieben wird.

  \item \textbf{Einsatz in globalen Standards:}
    ECC-Kurven wie \texttt{P-256} (NIST), \texttt{Curve25519} (Bernstein 2006)
    oder \texttt{secp256k1} (Bitcoin, Ethereum) sichern
    TLS/HTTPS-Verbin\-dungen, Smartphone-Verschlüsselung, Blockchain-Transaktionen
    und digitale Signaturen (ECDSA, EdDSA).
    Ein typischer 256-Bit-ECC-Schlüssel bietet vergleichbare Sicherheit wie ein
    3072-Bit-RSA-Schlüssel, bei dramatisch besserem Rechenaufwand.

  \item \textbf{Schnittstelle zu Faltings:}
    Die Wahl sicherer Kurven erfordert tiefes Verständnis der Galoisdarstellungen
    $\rho_\ell: G_K \to GL_{2g}(\Z_\ell)$. Das Irreduzibilitätskriterium
    (Satz~\ref{thm:irred_krit}) und die Analyse des Frobenius-charakteristischen Polynoms
    $P_{J,\mathfrak p}(T)$ sind direkte Werkzeuge bei der Überprüfung der
    Kryptosicherheit: Ein reduzibler Tate-Modul würde eine Angriffsfläche bieten.

  \item \textbf{Isogenie-basierte Kryptographie:}
    Die neueste Generation post-quantensicherer Protokolle nutzt \emph{Isogenien}
    zwischen elliptischen Kurven als Einwegfunktionen. Der Algorithmus
    \texttt{SIDH} (Supersingular Isogeny Diffie-Hellman, Jao-De Feo 2011),
    die Variante \texttt{CSIDH} (Castryck et al.\ 2018) und der NIST-Kandidat
    \texttt{SQIsign} (De Feo et al.\ 2020) basieren alle auf der Struktur
    der Isogeniegraphen, die direkt durch Faltings' Isogenie-Endlichkeitssätze
    kontrolliert werden. Der 2022 erfolgte Angriff auf SIDH (Castryck-Decru)
    verdeutlichte, wie tiefes algebraisch-geometrisches Wissen — nämlich über
    Torsionspunkte auf höher-dimensionalen abelschen Varietäten — für die
    Kryptanalyse entscheidend ist.
\end{itemize}

\subsubsection{Fehlerkorrigierende Codes}

\begin{itemize}[leftmargin=2em,itemsep=4pt]
  \item \textbf{Algebraisch-geometrische Codes (AG-Codes):}
    Goppa (1981) und Tsfasman-Vlăduţ-Zink (1982) zeigten, dass algebraische Kurven
    mit vielen $\F_q$-rationalen Punkten optimale fehlerkorrigierende Codes liefern.
    Die \emph{Hamming-Schranke} und \emph{Singleton-Schranke} werden durch
    AG-Codes übertroffen, sobald die Kurven vom Geschlecht $g \geq 1$ genügend
    viele rationale Punkte besitzen.

  \item \textbf{Verbindung zu Faltings:}
    Das Weil-Theorem für Kurven (Weil 1948) zeigt: Eine Kurve $C$ vom Geschlecht $g$
    über $\F_q$ erfüllt $|C(\F_q) - (q+1)| \leq 2g\sqrt{q}$ (\emph{Hasse-Weil-Schranke}).
    Das Zusammenspiel dieser Schranke mit dem Faltings-Theorem über Zahlkörper
    zeigt, dass hochgratige Kurven (die über $\Q$ wenige Punkte haben)
    über endlichen Körpern hochoptimale Code-Parameter erzeugen können.

  \item \textbf{Anwendungen in Technik und Telekommunikation:}
    AG-Codes auf Hermite-Kurven und Suzuki-Kurven werden in der
    Raumfahrtkommunikation (NASA Deep Space Network), in optischen Netzwerken
    (Reed-Solomon als Spezialfall), in RAID-Systemen sowie in Quantum-Error-Correction-Codes
    (QECC) für Quantencomputer eingesetzt.
\end{itemize}

%% ------------------------------------------------------------
\subsection{Computeralgebra, Algorithmen und Zahlentheorie}
%% ------------------------------------------------------------

\begin{itemize}[leftmargin=2em,itemsep=4pt]
  \item \textbf{Berechenbare Arithmetik elliptischer Kurven:}
    Effiziente Algorithmen zur Berechnung von $|E(\F_q)|$ — etwa der
    \emph{Schoof-Algorithmus} (1985) und seine Verbesserungen durch
    Atkin und Elkies (SEA-Algorithmus) — nutzen genau die $\ell$-adischen
    Tate-Moduln aus Abschnitt~\ref{sec:tate} und berechnen das charakteristische
    Polynom des Frobenius. Dies ist der direkteste algorithmische Ausläufer
    von Faltings' Theorie.

  \item \textbf{Chabauty-Coleman-Methode:}
    Die klassische Methode von Chabauty (1941) und Coleman (1985) berechnet
    $C(\Q)$ für $\rank J(\Q) < g$ durch $p$-adische Integration auf $C$.
    Balakrishnan und Dogra (2017–2021) haben durch Einbezug der
    nicht-abelschen Chabauty-Methode die Schranke auf $\rank J(\Q) < 2g-2$
    erweitert und konkrete Algorithmen in \texttt{SageMath} und \texttt{Magma}
    implementiert. Diese Algorithmen berechnen heute routinemäßig alle rationalen
    Punkte auf Kurven vom Geschlecht 2 und 3.

  \item \textbf{Das \texttt{LMFDB}-Projekt:}
    Die \emph{L-Functions and Modular Forms Database} (\texttt{lmfdb.org}) katalogisiert
    tausende von elliptischen Kurven, höhergenusigen Kurven, abelschen Varietäten
    und zugehörigen Galoisdarstellungen über Zahlkörpern. Die Datenbank speichert
    explizit die Frobenius-Polynome $P_{J,\mathfrak p}(T)$ und ermöglicht
    maschinenlesbare Überprüfungen des Irreduzibilitätskriteriums
    (Satz~\ref{thm:irred_krit}) in großem Maßstab.

  \item \textbf{Formale Verifikation:}
    Mit dem Aufkommen formaler Beweissysteme (\texttt{Lean 4}, \texttt{Coq},
    \texttt{Isabelle/HOL}) beginnt die Community, grundlegende Sätze wie den
    von Mordell-Weil formell zu verifizieren. Das \emph{Lean4-Mathlib}-Projekt
    (2023) hat Teile der Theorie elliptischer Kurven formalisiert.
    Die Formalisierung des Satzes von Faltings selbst ist ein Langzeitziel,
    das die Entwicklung neuer Abstraktionen in der formalen Zahlentheorie
    erfordern wird.
\end{itemize}

%% ------------------------------------------------------------
\subsection{Anwendungen in anderen Wissenschaften}
%% ------------------------------------------------------------

\subsubsection{Quanteninformatik und Quantenkryptographie}

\begin{itemize}[leftmargin=2em,itemsep=4pt]
  \item \textbf{Post-Quanten-Kryptographie:}
    Quantencomputer bedrohen ECC durch Shors Algorithmus (1994), der ECDLP
    in polynomialer Zeit löst. Als Reaktion entwickelt das NIST-Post-Quantum-Standardisierungsprojekt
    (abgeschlossen 2024) neue kryptographische Standards. Neben gitterbasierten
    Systemen (\texttt{CRYSTALS-Kyber}, \texttt{CRYSTALS-Dili\newline thium})
    sind isogenie-basierte Systeme wie \texttt{SQIsign} direkte Nachfolger
    der durch Faltings inspirierten Isogenienkryptographie.

  \item \textbf{Quantenfehlerkorrektur:}
    Topologische Quantencodes (Kitaev Toric Code, Surface Codes) basieren auf
    homologischen Eigenschaften von Mannigfaltigkeiten. Die Verbindung zu
    AG-Codes und damit zu algebraischen Kurven wird in einem aktiven Forschungsgebiet
    untersucht: \emph{Quantum LDPC Codes} aus algebraisch-geometrischen Codes
    versprechen optimale Parameter für Quantencomputer.
\end{itemize}

\subsubsection{Physik der Kondensation und Statistik}

\begin{itemize}[leftmargin=2em,itemsep=4pt]
  \item \textbf{Zufällige Matrizen und $L$-Funktionen:}
    Die \emph{Riemann-Vermutung} und ihre Verallgemeinerungen (GRH über Zahlkörpern,
    Weil-Vermutungen für Kurven) verbinden die Nullstellen von $L$-Funktionen mit
    dem Spektrum zufälliger Matrizen (GUE, GOE). Die Frobenius-Eigenwerte
    $\alpha_p, \bar\alpha_p$ der Hasse-Weil-Zetafunktion einer elliptischen Kurve
    zeigen genau die statistischen Eigenschaften des GUE-Ensembles — eine
    fundamentale Brücke zwischen Zahlentheorie und mathematischer Physik.

  \item \textbf{Gittersysteme in der Kondensationstheorie:}
    Die Néron-Tate-Höhe definiert eine positiv definite quadratische Form auf
    $J(K) \otimes \R$, die strukturell identisch mit der Energie eines
    physikalischen \emph{Coulomb-Gitters} ist. Die Verteilung rationaler Punkte
    minimaler Höhe entspricht der Grundzustandsverteilung eines Gittermodells,
    eine Verbindung, die in der arithmetischen Dynamik (Silverman, Call) systematisch
    untersucht wird.
\end{itemize}

\subsubsection{Biologische Mathematik und Netzwerktheorie}

\begin{itemize}[leftmargin=2em,itemsep=4pt]
  \item \textbf{Phylogenetische Bäume und tropische Geometrie:}
    Tropische Kurven, die als metrische Graphen auftreten, modellieren
    \emph{phylogenetische Bäume} in der Evolutionsbiologie.
    Die Methoden der tropischen Geometrie — direkt aus Faltings' Theorie
    der stabilen Kurvenentartungen hergeleitet — wurden von Speyer-Sturmfels (2004)
    zur Analyse von DNA-Sequenzdaten eingesetzt.

  \item \textbf{Topologische Datenanalyse (TDA):}
    Die \emph{persistente Homologie} (Edelsbrunner, Carlsson) misst geometrische
    Strukturen in Datenwolken durch topologische Invarianten. Methoden aus der
    algebraischen Geometrie, insbesondere Garbentheorie und Kohomologie,
    bilden die theoretische Grundlage moderner TDA-Algorithmen.
    AG-Codes werden in der TDA als Fehlerkorrekturschicht in der
    robusten Datenanalyse eingesetzt.
\end{itemize}

%% ------------------------------------------------------------
\subsection{Industrie und Technologietransfer}
%% ------------------------------------------------------------

\subsubsection{Finanzsektor und Blockchain}

\begin{itemize}[leftmargin=2em,itemsep=4pt]
  \item \textbf{Blockchain-Kryptographie:}
    Nahezu alle großen Kryptowährungen — Bitcoin, Ethereum, Solana —
    verwenden elliptische Kurven für digitale Signaturen.
    Bitcoin nutzt \texttt{secp256k1}: $y^2 = x^3 + 7$ über $\F_p$ mit einer
    256-Bit-Primzahl $p$. Ethereum-2.0 verwendet \texttt{BLS12-381}, eine
    \emph{Paarungskurve} (Barreto-Lynn-Scott), die auf abelschen Varietäten
    der Dimension 2 (Jacobische einer Genus-2-Kurve) basiert und effiziente
    Threshold-Signaturen ermöglicht. Das wirtschaftliche Volumen dieser Systeme
    beläuft sich auf über 2 Billionen US-Dollar (Stand 2026).

  \item \textbf{Zero-Knowledge-Proofs:}
    Moderne Zero-Knowledge-Beweissysteme (ZK-SNARKs, ZK-STARKs) für
    \emph{vertrauliche Transaktionen} (Zcash, StarkNet) verwenden
    paarungsbasierte Kryptographie und bivariate Polynominterpolation
    über elliptischen Kurven. Groth16, PlonK und andere Proof-Systeme
    basieren auf dem \emph{Pairing-Bilinearity}, das durch die Weil-Paarung
    auf abelschen Varietäten — direkt aus dem Faltings-Rahmen — bereitgestellt wird.

  \item \textbf{Finanzielle Modellierung:}
    Höhenfunktionen und $L$-Funktionen algebraischer Varietäten wurden als
    Modelle für Zinsstrukturkurven und Optionspreisformeln untersucht.
    Insbesondere die Verbindung zwischen den Eigenwerten des Frobeniusoperators
    und Wahrscheinlichkeitsmaßen auf dem Spektrum führt zu randständigen,
    aber potenziell fruchtbaren Querschnittsideen zwischen Finanzmathematik
    und Zahlentheorie.
\end{itemize}

\subsubsection{Telekommunikation und Hardwaredesign}

\begin{itemize}[leftmargin=2em,itemsep=4pt]
  \item \textbf{Implementierung in Hardware:}
    ECC-Operationen auf \texttt{Curve25519} und \texttt{P-256} sind in
    allen modernen CPUs (ARM TrustZone, Intel SGX), Sicherheitschips
    (TPM 2.0, Secure Enclave von Apple) und Smart-Cards direkt in Hardware implementiert.
    Optimierungen der Skalarmultiplikation $P \mapsto nP$ exploitieren die
    $\Z_\ell$-Modulstruktur des Tate-Moduls zur effizienten Berechnung.

  \item \textbf{Quantenkommunikation:}
    Quantenschlüsselverteilung (QKD) nach BB84 und weiterentwickelte Protokolle
    verwenden in der klassischen Authentifizierungsschicht ECC-basierte Signaturen.
    Der Ãœbergang zu post-quantensicheren Protokollen (isogenie-basiert oder
    gitterbasiert) ist ein aktives Standardisierungsfeld (ETSI, ISO/IEC~18033).

  \item \textbf{Autonome Systeme und PKI:}
    Public-Key-Infrastrukturen (PKI) für eingebettete Systeme, Fahrzeugkommunikation
    (V2X, \texttt{IEEE~1609.2}) und Industriesteuerungen (IEC~62351)
    setzen auf ECC, da der geringe Schlüssel- und Signaturumfang
    eng begrenzte Ressourcen schont.
\end{itemize}

\subsubsection{Künstliche Intelligenz und Maschinelles Lernen}

\begin{itemize}[leftmargin=2em,itemsep=4pt]
  \item \textbf{Homomorphe Verschlüsselung:}
    \emph{Fully Homomorphic Encryption} (FHE, Gentry 2009) ermöglicht Berechnungen
    auf verschlüsselten Daten — dies ist die Grundlage vertraulichen maschinellen
    Lernens (\emph{privacy-preserving ML}).
    Aktuelle FHE-Schemata wie CKKS, BGV, BFV beruhen auf
    \emph{Ring-LWE} über ausgewählten Polynomringen, deren Parameterauswahl
    durch zahlentheoretische Methoden gesteuert wird.
    Gitterbasierte FHE-Parameter werden durch die Geometrie von
    Gitterreduktionsalgorithmen (LLL, BKZ) optimiert — Algorithmen,
    die auf derselben arithmetischen Geometrie basieren wie die
    Faltings-Höhentheorie.

  \item \textbf{Sichere mehrparteiige Berechnung (MPC):}
    MPC-Protokolle (Yao, GMW, SPDZ) für verteiltes maschinelles Lernen
    verwenden geheimes Teilen über endlichen Körpern $\F_p$,
    wobei Shamir's Secret Sharing auf polynomialer Interpolation auf
    algebraischen Kurven (Hermite-Kurven) basiert.
    In großskaligen Federated-Learning-Architekturen sind solche AG-Code-basierten
    MPC-Protokolle den klassischen Alternativen überlegen.
\end{itemize}

%% ------------------------------------------------------------
\subsection{Aktuelle Forschungsfronten und offene Probleme}
%% ------------------------------------------------------------

Wir schließen mit einem Überblick über die wichtigsten offenen Fragen, die direkt
aus dem Satz von Faltings und seinen Methoden hervorgehen.

\begin{enumerate}[leftmargin=2em,itemsep=6pt, label=\textbf{(\arabic*)}]
  \item \textbf{Effektiver Satz von Faltings.}
    Gibt es einen Algorithmus, der — gegeben $C/K$ mit $g(C) \geq 2$ —
    alle Punkte in $C(K)$ berechnet? Bisher ist kein solcher Algorithmus bekannt.
    Kims nicht-abelsche Chabauty-Methode liefert in Spezialfällen effektive Ergebnisse,
    eine allgemeine Lösung fehlt. Dies ist wohl das wichtigste offene Problem
    der algorithmischen Zahlentheorie.

  \item \textbf{Gleichmäßige Schranken für $|C(K)|$.}
    Der Uniform-Mordell-Lang-Satz (Dimitrov-Gao-Habegger 2021) liefert eine
    Schranke $|C(K)| \leq C(g, [K:\Q])$, aber die genaue Abhängigkeit von $g$
    ist noch ungeklärt. Gilt $|C(\Q)| \leq 2g - 2 + \varepsilon$ für alle
    $C/\Q$ mit $g \geq 2$ außer endlich vielen? Eine solche scharfe uniforme Schranke
    ist offen.

  \item \textbf{Fontaine-Mazur-Vermutung.}
    Die Fontaine-Mazur-Vermutung besagt: Jede geometrische, irreduzible
    $\ell$-adische Galoisdarstellung $\rho: G_\Q \to GL_n(\Q_\ell)$
    stammt von einer algebraischen Varietät über $\Q$. Satz~\ref{thm:galois_struktur}
    liefert Belege im Fall $n = 2g$, aber die allgemeine Vermutung ist offen.

  \item \textbf{Geometrisches Langlands-Programm (vollständige Quantisierung).}
    Die jüngste Arbeit (2024) bewies das geometrische Langlands-Programm in der
    klassischen (de-Rham-)Version. Die $p$-adische und quantisierte Version
    (geometrisches $p$-adisches Langlands) ist ein aktives Forschungsfeld, in dem
    Faltings' perfectoide Methoden eine zentrale Rolle spielen.

  \item \textbf{Böcherer-Vermutung und Modulformen.}
    Böcherers Vermutung (1985) verbindet Fourier-Koeffizienten von
    Siegel-Modulformen der Dimension $g$ mit speziellen Werten von
    $L$-Funktionen. Sie ist ein Analogon des Modularsatzes (Taylor-Wiles)
    für höherdimensionale abelsche Varietäten. Ein Beweis würde tiefgreifende
    neue Anwendungen der Faltings-Theorie erfordern.

  \item \textbf{Riemann-Hypothese für $L$-Funktionen.}
    Die Riemann-Vermutung über die Lage der Nullstellen von $\zeta(s)$
    ist äquivalent zu Aussagen über die Verteilung von Primzahlen,
    die in geeigneter Sprache als Aussagen über die Frobenius-Eigenwerte
    in der Faltings-Theorie reformuliert werden können.
    Weil's Theorem für Kurven (1948) und Deligne's Theorem für Varietäten (1974)
    — beide enge Verwandte der Faltings-Methodik — lösten analoge Vermutungen
    für endliche Körper. Der Fall des klassischen $\zeta(s)$ bleibt offen.

  \item \textbf{Post-Quanten-Sicherheitsanalyse isogenie-basierter Systeme.}
    Nach dem Cast\-ryck-Decru-Angriff auf SIDH (2022) ist die Sicherheit
    isogenie-basierter Kryptosysteme neu zu bewerten. Insbesondere ist unklar,
    ob das CSIDH-System Quantensicherheit bietet. Die Antwort liegt in der
    Geometrie von Isogeniegraphen supersingularer elliptischer Kurven —
    ein direkter Ausläufer der Faltings-Theorie abelscher Varietäten.
\end{enumerate}

%% ============================================================
%% Literatur
%% ============================================================
\begin{thebibliography}{99}
\bibitem{faltings83}
G.\ Faltings,
\textit{Endlichkeitssätze für abelsche Varietäten über Zahlkörpern},
Inventiones Mathematicae \textbf{73} (1983), 349--366.

\bibitem{mordell22}
L.J.\ Mordell,
\textit{On the Rational Solutions of the Indeterminate Equations of the Third and Fourth Degree},
Proc.\ Cambridge Phil.\ Soc.\ \textbf{21} (1922), 179--192.

\bibitem{silverman86}
J.H.\ Silverman,
\textit{The Arithmetic of Elliptic Curves},
Graduate Texts in Mathematics \textbf{106}, Springer, New York, 1986.

\bibitem{silverman94}
J.H.\ Silverman,
\textit{Advanced Topics in the Arithmetic of Elliptic Curves},
Graduate Texts in Mathematics \textbf{151}, Springer, New York, 1994.

\bibitem{hindry00}
M.\ Hindry, J.H.\ Silverman,
\textit{Diophantine Geometry: An Introduction},
Graduate Texts in Mathematics \textbf{201}, Springer, New York, 2000.

\bibitem{bost96}
J.-B.\ Bost,
\textit{Périodes et isogenies des variétés abéliennes sur les corps de nombres},
Séminaire Bourbaki 1994/95, Astérisque \textbf{237} (1996), Exp.\ No.\ 795.

\bibitem{parshin68}
A.N.\ Parshin,
\textit{Algebraic curves over function fields I},
Math.\ USSR Izvestija \textbf{2} (1968), 1145--1170.

\bibitem{bombieri90}
E.\ Bombieri,
\textit{The Mordell conjecture revisited},
Ann.\ Scuola Norm.\ Sup.\ Pisa \textbf{17} (1990), 615--640.

\bibitem{dimitrov21}
V.\ Dimitrov, Z.\ Gao, P.\ Habegger,
\textit{Uniformity in Mordell-Lang for curves},
Annals of Mathematics \textbf{194} (2021), 237--298.

\bibitem{tate66}
J.\ Tate,
\textit{Endomorphisms of abelian varieties over finite fields},
Inventiones Mathematicae \textbf{2} (1966), 134--144.

\bibitem{wiles95}
A.\ Wiles,
\textit{Modular elliptic curves and Fermat's Last Theorem},
Annals of Mathematics \textbf{141} (1995), 443--551.

\bibitem{hriljac85}
P.\ Hriljac,
\textit{Heights and Arakelov's intersection theory},
Pacific J.\ Math.\ \textbf{87} (1985), 359--378.
\end{thebibliography}

\end{document}

